\chapter{電磁気・物性}
\section{電荷・場・回路をつなぐ}
電磁気の問題では、電荷と電流が場を作る過程、場が電荷へ力を及ぼす過程、
場に蓄えられるエネルギーを分けて整理する。
回路は、導線と素子へ分布する場を端子電圧と電流で近似したモデルである。
電荷一定か電圧一定か、静電平衡か時間変化する場か、直列か並列かを最初に確定する。
この共通本文は静電場・容量・直流・磁場・誘導・RLCの27問に対応する初稿である。
磁性・物性等の未執筆問題も一覧へ保持し、章全体の完成とはしない。

\section{点電荷と電場}\label{theory:em:coulomb}
真空中の点電荷間の力の大きさは
\[
 F=\frac{|q_1q_2|}{4\pi\varepsilon_0r^2},\qquad \mathbf F=q\mathbf E
\]
である。第一式は二つの電荷が作る相互作用、第二式の$\mathbf E$は対象電荷以外が作る場である。
正電荷の力は電場と同じ方向、負電荷では逆方向になる。
数値問題では電荷をC、距離をmへそろえ、最後にNになることを確かめる。
同符号か異符号かは力の向きを決め、大きさを求める絶対値とは別に扱う。\autocite{mit_em02}

\section{電位差と外力の仕事}\label{theory:em:potential}
静電場では$\mathbf E=-\nabla V$であり、PからQへの移動で
\[
 V_Q-V_P=-\int_P^Q\mathbf E\cdot d\boldsymbol\ell,
 \quad \Delta U=q(V_Q-V_P),\quad W_{\rm field}=-\Delta U
\]
となる。ゆっくり動かし、運動エネルギー変化を無視できるとき、外力の仕事は$W_{\rm ext}=\Delta U$である。
電場方向へ負電荷を動かす場合、電位は下がるが$q<0$なので位置エネルギーは増える。
したがって正の外力仕事と負の電位変化は矛盾しない。
候補が正の電位差だけを示す場合には、大小差$|V_Q-V_P|$と、向きを定めた差を本文で明示する。
時間変化する磁束のある回路では誘導電場が生じるため、静電位だけで全起電力を表せるとは限らない。\autocite{mit_em03}

\section{ガウス法則と導体}\label{theory:em:gauss}
真空中では$\int_S\mathbf E\cdot\mathbf n\,dS=Q/\varepsilon_0$である。
一様帯電平面をまたぐ薄い柱を取ると、側面の寄与は0、両端面の寄与は各$EA$となるため
$2EA=|\sigma|A/\varepsilon_0$、すなわち$E=|\sigma|/(2\varepsilon_0)$を得る。
一枚の平面と、逆符号二枚の間で場を足した場合の係数を区別する。
静電平衡の導体内部は$E=0$である。導体内部に場があれば自由電荷の移動が続くためである。
孤立した球の外側にある逆二乗場を、導体内部へ延長しない。\autocite{mit_em04}

\section{容量と接続}\label{theory:em:capacitor}
容量は$C=Q/V$で、理想平行平板では$C=\varepsilon A/d$となる。
容量の単位FはC/Vである。$Q$は片方の極板の電荷の大きさを表す。
未充電のコンデンサーを直列接続すると、孤立した中間導体の正味電荷が0のまま保存されるため、
各素子の電気量の大きさは等しくなる。その電圧を足すと$1/C_{\rm eq}=\sum 1/C_i$である。
並列では電圧が共通なので電荷を足し、$C_{\rm eq}=\sum C_i$となる。
抵抗の直列・並列とは逆の合成関係に見えるが、共通量と加算量から導けば暗記を必要としない。\autocite{mit_em05}

\section{電場エネルギー}\label{theory:em:energy}
固定形状の線形コンデンサーに微小電荷$dq$を運ぶ仕事は$dW=Vdq=(q/C)dq$である。
0からQまで積分して
\[
 U=\int_0^Q\frac qC\,dq=\frac{Q^2}{2C}=\frac12CV^2
\]
を得る。$QV$では充電中の電圧変化を無視しており、2倍になる。
エネルギーの数値計算で$\mu$を$10^{-6}$へ変換し、電圧の二乗に注意する。
LC回路で用いる$LI^2/2$も場へ蓄えるエネルギーであり、抵抗で熱になる消費エネルギーとは違う。

\section{分極・強誘電性・損失}\label{theory:em:dielectric}
分極$\mathbf P$は単位体積あたりの電気双極子モーメントで、
$\mathbf D=\varepsilon_0\mathbf E+\mathbf P$と表す。
通常の誘電体で外場に応じて分極することと、強誘電体の自発分極を区別する。
強誘電性は外部電場なしでも生じる分極と、電場によるその反転を特徴とする。
交流で分極の応答に遅れや散逸があると、供給エネルギーの一部が熱となる誘電損失が生じる。
電解コンデンサーの酸化アルミニウム・酸化タンタル膜は薄い絶縁層として使われる。
誘電率が大きいことだけで、損失や耐電圧もよいとは判断しない。\autocite{mit_em05,mit_ferroelectric}

\section{固定電荷での極板間引力}\label{theory:em:force}
端効果のない一様な線形誘電体を仮定する。Q一定で間隔$d$を変えると
\[
 U(Q,d)=\frac{Q^2d}{2\varepsilon A},\qquad
 F_{\rm attraction}=\left(\frac{\partial U}{\partial d}\right)_Q
 =\frac{Q^2}{2\varepsilon A}
\]
である。これは引力の大きさであり、間隔を増やす方向の力成分は負になる。
Qが一定なら間隔の変更だけでは引力は変わらないが、誘電率を5倍にすれば引力は1/5となる。
容量が変更前と同じでも、間隔微分が同じとは限らない。
電圧一定の場合は電源とのエネルギーの授受があるため、$CV^2/2$を固定Qの式と同じように微分しない。
固定Vでの大きさは$\varepsilon AV^2/(2d^2)$である。条件を先に選ぶ必要がある。

\section{直流と抵抗}\label{theory:em:dc}
電流は断面を通る符号付き電荷の時間変化率$I=dQ/dt$で、1 A=1 C/sである。
負の担い手の移動方向と電流方向は逆になる。電子に限らず、イオンも電流を運べる。
オーム性導体では$V=IR$、電力は$P=VI=I^2R=V^2/R$となる。
抵抗を並列につなぐと電圧が共通で$I=\sum V/R_i$より$1/R_{\rm eq}=\sum1/R_i$を得る。
合成抵抗が最小の構成抵抗より小さいことも検算になる。\autocite{mit_em06,mit_em07}

\section{抵抗率と担い手}\label{theory:em:resistivity}
均一な棒の抵抗は$R=\rho L/A$で、材料の抵抗率と形状の寄与を分離する。
金属では格子振動、不純物、欠陥等による散乱が抵抗へ寄与する。
室温近傍の変化は$\rho(T)\simeq\rho(T_0)[1+\alpha(T-T_0)]$で近似することが多いが、全温度域の厳密な法則ではない。
真性半導体では温度上昇で電子と正孔が熱励起され、担い手の増加が抵抗減少の主因となる。
低温の金属での特定のべき則を、多くの導体についての一般則へ広げない。\autocite{mit_em06,mit_semiconductor}

\section{電流による場と荷電粒子の運動}\label{theory:em:biot}\label{theory:em:lorentz}
定常電流素片が作る磁束密度はビオ・サバールの法則
$d\mathbf B=\mu_0I(d\boldsymbol\ell\times\hat{\mathbf r})/(4\pi r^2)$で与えられる。
直線電流の周囲の場は円周方向へ向き、局所的に電流方向へ垂直になる。
一方、場の中で動く粒子への磁気力は$q\mathbf v\times\mathbf B$である。
その仕事率は$\mathbf F\cdot\mathbf v=0$なので速さを変えず、方向を変える。
一様Bへ垂直な速度で入り、他の力がなければ円運動となり、半径は$mv/(|q|B)$である。
場を作る法則と、その場から力を受ける法則を混同しない。\autocite{mit_em09,mit_lorentz}

\section{ソレノイドと透磁率}\label{theory:em:solenoid}
長いコイルの端部から離れた内部では、巻数密度$n=N/\ell$を使い$B=\mu nI$と近似する。
全巻数Nをそのまま掛けると長さの次元が足りない。
磁性体があるときは、出題で指定された線形範囲の$\mu=\mu_0\mu_r$を用いる。
有限長の端部の場、材料の飽和や履歴はこの簡単な式の外にある。
TからmTへの換算では$10^3$を掛ける。\autocite{mit_em09}

\section{誘導とマクスウェル方程式}\label{theory:em:induction}\label{theory:em:maxwell}
一巻きの磁束を$\Phi_B$とすると$\mathcal E=-N\,d\Phi_B/dt$である。
負号は変化を妨げる向きであり、大きさの問題では$N|d\Phi_B/dt|$を使う。
真電荷を自由電荷として扱う巨視的な表記では
\begin{align*}
 \nabla\cdot\mathbf D&=\rho_{\rm free},&\nabla\cdot\mathbf B&=0,\\
 \nabla\times\mathbf E&=-\frac{\partial\mathbf B}{\partial t},&
 \nabla\times\mathbf H&=\mathbf i+\frac{\partial\mathbf D}{\partial t}.
\end{align*}
磁束変化から電場を得る法則と、電束密度変化から変位電流を得る項を区別する。
$\nabla\cdot\mathbf B=0$は、通常の電磁気学で磁気単極子を源としないことを表す。
電荷保存は$\partial_t\rho+\nabla\cdot\mathbf i=0$で、流出を正とした総電荷変化には負号が付く。
「出入り」を流入と流出の単なる和として扱わない。\autocite{mit_em10,mit_em13}

\section{LC振動と直列共振}\label{theory:em:lc}\label{theory:em:rlc}
無抵抗LC回路では$Q^2/(2C)+LI^2/2$が保存する。
最大電流時はQ=0となり、初期にコンデンサーへ蓄えたエネルギーが全てコイルへ移る。
交流電源で駆動する理想直列RLC回路では
\[
 Z=R+j\left(\omega L-\frac1{\omega C}\right),\quad
 |Z|^2=R^2+\left(\omega L-\frac1{\omega C}\right)^2
\]
となる。括弧が0となる$f_0=1/(2\pi\sqrt{LC})$で$|Z|$は最小Rになる。
抵抗で減衰する自由振動の周波数と、交流定常状態の電流共振を混同しない。
L,CをSIへ直してから計算し、最後にHzをkHzへ換算する。\autocite{mit_em11,mit_em12}

\section{場のエネルギー流}\label{theory:em:poynting}
電磁場のエネルギー流密度は$\mathbf S=\mathbf E\times\mathbf H$、単位はW/m$^2$である。
外積なのでE,H双方へ垂直であり、流れる方向と単位面積あたりの仕事率を同時に表す。
力の向きや電荷の移動方向と同一視しない。
電磁波だけでなく、回路へエネルギーが運ばれる記述にも使われる。\autocite{mit_em13}

\section{磁化を符号・秩序・履歴に分ける}\label{theory:em:magnetism}
単位体積あたりの磁気モーメントを磁化\(\mathbf M\)とする。SIでは
\(\mathbf B=\mu_0(\mathbf H+\mathbf M)\)。等方的な弱い線形応答では
\(\mathbf M=\chi\mathbf H\)で、反磁性は負の\(\chi\)、常磁性は小さな正の\(\chi\)を持つ。
強磁性は外場なしの自発磁化を持つ秩序であり、多磁区試料の正味磁化が小さいことと両立する。
反強磁性は典型的には反平行の副格子磁化が打ち消し合う秩序で、Néel温度を超えると秩序を失う。
通常の強磁性の秩序を失う温度はCurie温度と呼ぶ。残留磁化は外場を取り去った後の履歴による磁化で、
磁区配置と保磁力を含む材料の状態に依存する。\cite{openstax_magnetism}

室温の代表的強磁性元素はFe・Co・Niである。過去問のMnをCoに読み替えない。
磁性は温度・結晶相・組成にも依存する。中性子回折は核による構造散乱だけでなく磁気散乱にも使える。
Slater--Pauling曲線は合金の電子数と磁気モーメント、Ising模型はスピン相互作用と秩序を扱う。
一方、Wiedemann--Franz則の\(\kappa_e/(\sigma T)\)は電子の熱電輸送係数の比である。
磁性に関する相対比較で区別し、磁場やスピンの影響を全く受けないと断定しない。\cite{em51_mit_magnetic_readings,em51_wf}

\section{局在モーメントと伝導電子}\label{theory:em:pauli}
独立した局在モーメント、弱い外場、高温側のモデルでは、磁化の方向が熱揺らぎで乱され、
磁化率はCurie則\(\chi=C/T\)に従う。Tは絶対温度である。
相互作用を含むCurie--Weiss近似の\(C/(T-\theta)\)も、転移点近傍や全温度の厳密則ではない。
「常磁性なら必ず同じ温度依存」という暗記を避ける。

Pauli常磁性は金属の伝導電子のスピンへの応答で、占有のFermi統計が重要になる。
低温では全電子が独立に方向を変える局在モーメントのモデルを使えず、Fermi準位近傍の状態が応答する。
原子核まわりの軌道運動を起源とするという過去問の説明はこの機構と違う。
自由電子の軌道応答であるLandau反磁性、束縛電子の応答を分ける。\cite{em51_pauli}

超伝導のMeissner状態では磁束が排除されるが、温度・磁場の条件がある。
II型の混合状態では量子化磁束が侵入するため、「超伝導の全状態で完全排除」とはしない。
電気抵抗ゼロという輸送性質だけから、磁化の全状態を決めない。

\section{形状で変わる反磁場と磁場勾配の力}\label{theory:em:demag}
試料の磁化は周囲と内部の場を変える。均一磁化した楕円体では、主軸方向ごとに
\[
 H_{\mathrm d,i}=-N_iM_i,\qquad N_x+N_y+N_z=1
\]
と書ける。NはSIでの無次元の反磁場係数で、形状に依存する。
一般の不均一磁化・任意形状を一つのスカラーNで表すとは限らない。\cite{em51_demag}
反磁場は磁化から生じる場であり、反磁性という材料応答の名前とは別である。

黒鉛にも反磁性があるため、磁場中で「全く力を受けない」という全称は誤りである。
弱い等方的な誘起磁化の小試料なら、背景場を変えない近似で
\(\mathbf F\simeq\chi V\nabla(B^2)/(2\mu_0)\)と考えられる。
反磁性では強い場から弱い場へ向かう。実際の黒鉛の異方性や試料形状は別に考える。
理想的な一様場の並進力ゼロと、非一様場で働き得る力を区別する。

\section{巻数比を磁束変化から導く}\label{theory:em:transformer}
理想変圧器の一次・二次が共通の磁束\(\Phi(t)\)を共有するなら、
誘導起電力の大きさはそれぞれ\(N_1|d\Phi/dt|\)、\(N_2|d\Phi/dt|\)である。
したがって同じ位相規約での電圧の大きさ、実効値又は最大値同士では
\[
 \frac{E_2}{E_1}=\frac{N_2}{N_1}.
\]
巻数が多い側ほど電圧が大きいことが検算となる。損失を無視した電力保存から
\(E_1I_1=E_2I_2\)なので電流比は逆になる。
磁束漏れ、巻線抵抗、飽和、周波数応答を含む実機の限界は、この理想比の条件の外である。\cite{em51_transformer_ref}

\section{誘導加熱と誘電加熱を損失で分ける}\label{theory:em:heating}
誘導加熱は時間変化する磁場で導電体に渦電流を作り、そのJoule損を利用する。
磁化履歴損は強磁性体等に寄与するが、すべての導電体に必須の損失ではない。
誘電加熱は交流電場に対する分極応答の損失を利用し、代表的な対象は誘電体である。
有限の導電率を持つ材料では導電損と分極損が共存する場合がある。
マイクロ波加熱も誘電損失を利用できるが、高周波加熱の全方式がGHz帯という意味ではない。\cite{em51_heating}

良導体の線形な半空間に入る調和場では、表面からの場・電流の振幅が
\(\exp(-z/\delta)\)の形で減衰する。代表的な表皮深さは
\[
 \delta=\sqrt{\frac{2}{\omega\mu\sigma}}
       =\sqrt{\frac{\rho}{\pi f\mu}} .
\]
深さ\(\delta\)で振幅が\(1/e\)になるため、内部への侵入が厳密にゼロではない。
抵抗率・透磁率・周波数の条件を読み、有限試料の全温度分布をこの式だけで決めない。
加熱された表面から内部への熱伝導は電磁場の侵入とは別の過程である。\cite{em51_heating,em51_skin}
プラズマも高周波場から加熱できる。融合プラズマのRF加熱、無電極放電トーチ、
固体の渦電流加熱の回路と粒子の吸収機構は同一ではない。\cite{em51_plasma}

\section{副殻の容量と半導体の四面体結合}\label{theory:em:electrons}
一個の空間軌道にはスピンが異なる最大2電子を収容する。
s・p・d・f副殻の空間軌道数は1・3・5・7で、総容量は2・6・10・14である。
原典の「軌道」の語を、一個の軌道と副殻群で区別する。
Ge原子の外殻配置\(4s^24p^2\)と、通常の結晶の4近接・四面体のsp3結合モデルを分ける。
sp2の三角平面構造をダイヤモンド型へ当てはめない。\cite{atomic_quantum,em51_ge,em51_ge_structure}

SiへBをドープするとアクセプターを介して正孔が主要な担い手になる。
金属での負電子の平均ドリフト、半導体の正孔、孤立原子のスピンと
バルク金属のPauli応答を区別する。\cite{mit_semiconductor,em51_pauli}

\section{圧電性の逆命題を使わない}\label{theory:em:ferroelectric}
圧電性は応力による電気応答と電場によるひずみの結合である。
強誘電性は外場なしの自発分極と、その向きの電場による反転を特徴とする。
石英は圧電性を持つが、通常の強誘電体のような反転する自発分極を持たない。
圧電性から強誘電性を保証する逆命題は成立しない。\cite{mit_ferroelectric,em51_quartz}
多ドメイン試料では局所の自発分極が存在しても、正味の分極は相殺され得る。

典型的な強誘電体ではCurie温度を越えて常誘電相へ移る。
一次相転移は秩序変数が不連続に変わる場合で、強誘電相転移では自発分極を用いて考える。
有限の測定電場での飽和値・履歴と、電場ゼロでの熱力学的な秩序変数を無条件に同一視しない。
原子の変位、配向の秩序、結晶対称性と分極をつなぎ、微視的な機構を全物質で一種類に固定しない。\cite{em51_ferro_transition}

\section{特性線・連続放射と法令上の電波}\label{theory:em:spectral}
真空中の電磁波では\(c=f\lambda\)、光子エネルギーは\(E=hf\)である。
赤外・可視・紫外の順に波長が短くなる。
X線管の放射は、電子の減速による連続的な制動放射に、内殻準位差に対応する特性線が重なる。
管全体のスペクトルが連続成分を含むことと、特性X線自体を連続スペクトルと呼ぶことは違う。\cite{em51_xray}
取得できた電波法第二条の周波数上限は300万MHz以下で、\(3\times10^{12}\)Hz以下へ換算される。
令和5年度の施行版の直接取得には失敗したため、その時点の条文照合は個別問題で未確定として残す。
用途に基づくマイクロ波の波長区分と、法令の電波の範囲を別々に読む。\cite{em51_radio_law}

\section{過去問での確認}
\section{静電場・容量}

\subsection{平成23年度 問8}
\label{q:h23-reg-applied-08}

\noindent\textbf{原典：} sources/H23/17\_applied\_science.pdf、PDF物理3頁。

\textbf{公式正答：} 3（single）。

\paragraph{問題の要旨} 充電したコンデンサーに蓄えられるエネルギーを計算する。

\Needspace{4\baselineskip}\paragraph{確認した条件・記述（独自要約）}\mbox{}\par

\noindent\textbf{容量：} 2.0 μF\par

\noindent\textbf{電圧：} 300 V\par

\paragraph{問われる核心} 容量に蓄える電場エネルギー

\paragraph{必要知識} 静電容量、μ=10\textsuperscript{-6}、電圧上昇の積分。

\paragraph{考え方・正答を選ぶ理由（査読待ち初稿）} \(U=CV^2/2=(2.0\times10^{-6})300^2/2=0.09\) Jで3を選ぶ。充電中の微小仕事Vdqを積分するため、最終電圧を全電荷へ掛けたQVそのものではない。

\Needspace{4\baselineskip}\paragraph{選択肢の検討（査読待ち初稿）}\mbox{}\par

\noindent\textbf{1：条件に不一致} \(9.0\times10^{-3}\) J。正しい値0.09 Jの1/10である。\par

\noindent\textbf{2：条件に不一致} \(1.8\times10^{-2}\) J。正しい値0.09 Jの1/5である。\par

\noindent\textbf{3：条件に一致} \(9.0\times10^{-2}\) J。容量を2.0×10\textsuperscript{-6} Fへ換算したCV\textsuperscript{2}/2に一致する。\par

\noindent\textbf{4：条件に不一致} \(1.8\times10^{-1}\) J。CV\textsuperscript{2}とするとこの値になるが、充電中の電圧上昇を積分する係数1/2がない。\par

\noindent\textbf{5：条件に不一致} \(9.0\times10^{-1}\) J。正しい値の10倍で、接頭語又は指数の扱いが違う。\par

\paragraph{再利用する考え方} 容量に蓄える電場エネルギーを、成立条件・単位・符号から再現する。電場と電束密度、磁場と磁束密度、電荷一定と電圧一定を読み分ける。

\paragraph{共通理論} \hyperref[theory:em:capacitor]{第\ref*{theory:em:capacitor}節}, \hyperref[theory:em:energy]{第\ref*{theory:em:energy}節}

\par 参考：\cite{mit_em05}。

\par 難易度 2（校準初稿）、優先度未評価。

\subsection{平成24年度 問6}
\label{q:h24-reg-applied-06}

\noindent\textbf{原典：} sources/H24/17\_applied\_science.pdf、PDF物理2頁。

\textbf{公式正答：} 3（single）。

\paragraph{問題の要旨} 真空中の2点電荷間の静電気力の大きさを求める。

\Needspace{4\baselineskip}\paragraph{確認した条件・記述（独自要約）}\mbox{}\par

\noindent\textbf{電荷：} +2.0×10\textsuperscript{-6} C、+1.5×10\textsuperscript{-6} C\par

\noindent\textbf{距離：} 0.30 m\par

\noindent\textbf{係数：} 9.0×10\textsuperscript{9} N·m\textsuperscript{2}/C\textsuperscript{2}\par

\paragraph{問われる核心} 点電荷のクーロン力

\paragraph{必要知識} 電荷の符号、逆二乗則、単位と接頭語。

\paragraph{考え方・正答を選ぶ理由（査読待ち初稿）} \(F=k|q_1q_2|/r^2=9.0\times10^9(2.0\times10^{-6})(1.5\times10^{-6})/0.30^2=0.30\) N、3である。同符号なので互いに離れる方向の反発力であり、ここでは大きさを問う。

\Needspace{4\baselineskip}\paragraph{選択肢の検討（査読待ち初稿）}\mbox{}\par

\noindent\textbf{1：条件に不一致} 0.10 N。距離の逆二乗で評価した0.30 Nに一致しない。\par

\noindent\textbf{2：条件に不一致} 0.15 N。0.30 Nの半分であり、クーロン力へ1/2を掛ける理由がない。\par

\noindent\textbf{3：条件に一致} 0.30 N。電荷の積と距離の逆二乗から0.30 Nとなる。\par

\noindent\textbf{4：条件に不一致} 1.0 N。与えられた電荷と距離では0.30 Nである。\par

\noindent\textbf{5：条件に不一致} 3.0 N。正しい値の10倍である。\par

\paragraph{再利用する考え方} 点電荷のクーロン力を、成立条件・単位・符号から再現する。電場と電束密度、磁場と磁束密度、電荷一定と電圧一定を読み分ける。

\paragraph{共通理論} \hyperref[theory:em:coulomb]{第\ref*{theory:em:coulomb}節}

\par 参考：\cite{mit_em02}。

\par 難易度 2（校準初稿）、優先度未評価。

\subsection{平成25年度 問3}
\label{q:h25-reg-applied-03}

\noindent\textbf{原典：} sources/H25/17\_applied\_science.pdf、PDF物理2頁。

\textbf{公式正答：} 2（single）。

\paragraph{問題の要旨} 静電場の閉曲面積分の法則名と、無限一様帯電平面の電場の大きさを選ぶ。

\Needspace{4\baselineskip}\paragraph{確認した条件・記述（独自要約）}\mbox{}\par

\noindent\textbf{式：} \(\int_S E_n\,dS=Q/\varepsilon_0\)\par

\noindent\textbf{平面：} 面電荷密度σ。両側に等しい大きさで逆向きの垂直電場。\par

\noindent\textbf{組合せ：} 法則名／電場の大きさ。\par

\paragraph{問われる核心} ガウス法則と対称性

\paragraph{必要知識} 電束、外向き法線、無限平面、二面の寄与。

\paragraph{考え方・正答を選ぶ理由（査読待ち初稿）} 平面をまたぐ薄い柱を閉曲面に選ぶ。側面の電場法線成分は0、二つの端面の寄与は各EAである。内部電荷の大きさは|σ|Aなので\(2EA=|\sigma|A/\varepsilon_0\)、2を選ぶ。一枚の帯電平面と、逆符号二枚の間の電場を混同しない。

\Needspace{4\baselineskip}\paragraph{選択肢の検討（査読待ち初稿）}\mbox{}\par

\noindent\textbf{1：条件に不一致} ガウス／\(|\sigma|/\varepsilon_0\)。法則名は合うが、両面から出る電束の係数2が抜ける。\par

\noindent\textbf{2：条件に一致} ガウス／\(|\sigma|/(2\varepsilon_0)\)。ガウス法則を薄い柱へ適用した2EA=|σ|A/ε0に一致する。\par

\noindent\textbf{3：条件に不一致} ガウス／\(2|\sigma|/\varepsilon_0\)。両面の電束を足すと電場は半分になり、2倍ではない。\par

\noindent\textbf{4：条件に不一致} クーロン／\(|\sigma|/\varepsilon_0\)。原典の閉曲面積分はガウス法則で、電場の値も2倍である。\par

\noindent\textbf{5：条件に不一致} クーロン／\(|\sigma|/(2\varepsilon_0)\)。電場値は合うが提示された法則名と違う。\par

\paragraph{再利用する考え方} ガウス法則と対称性を、成立条件・単位・符号から再現する。電場と電束密度、磁場と磁束密度、電荷一定と電圧一定を読み分ける。

\paragraph{共通理論} \hyperref[theory:em:gauss]{第\ref*{theory:em:gauss}節}

\par 参考：\cite{mit_em04}。

\par 難易度 2（校準初稿）、優先度未評価。

\subsection{平成26年度 問4}
\label{q:h26-reg-applied-04}

\noindent\textbf{原典：} sources/H26/17\_applied\_science.pdf、PDF物理2頁。

\textbf{公式正答：} 4（single）。

\paragraph{問題の要旨} 負電荷を電場方向へゆっくり動かす外力の仕事から、電位差と電場の大きさを選ぶ。

\Needspace{4\baselineskip}\paragraph{確認した条件・記述（独自要約）}\mbox{}\par

\noindent\textbf{電荷：} −2 C\par

\noindent\textbf{移動：} PからQ。電場の方向と移動方向は同じ。\par

\noindent\textbf{外力の仕事：} 80 J。ゆっくり移動し運動エネルギー変化を無視。\par

\noindent\textbf{距離：} 0.1 m\par

\noindent\textbf{候補の順：} 電位差／電場。電位差は候補に合わせ大きさを表す。\par

\paragraph{問われる核心} 電位・外力仕事と負電荷の符号

\paragraph{必要知識} 準静的移動、位置エネルギーqV、E=−grad V、V/m=\allowbreak{}N/C。

\paragraph{考え方・正答を選ぶ理由（査読待ち初稿）} 外力の仕事は\(W_{ext}=q(V_Q-V_P)\)で、80=\allowbreak{}(−2)ΔVよりΔV=−40 Vである。電場方向へ進むと電位が下がるが、負電荷の位置エネルギーは増える。候補の電位差の大きさは40 V、\(|E|=|\Delta V|/d\)=400 V/mなので4である。電場の仕事は−80 Jと逆符号になる。

\Needspace{4\baselineskip}\paragraph{選択肢の検討（査読待ち初稿）}\mbox{}\par

\noindent\textbf{1：条件に不一致} 160 V／320 V·C。外力80 Jを電荷の絶対値2 Cで割ると40 Vであり、160 Vではない。\par

\noindent\textbf{2：条件に不一致} 160 V／1600 V/m。電位差の大きさは40 Vで、電場も400 V/mである。\par

\noindent\textbf{3：条件に不一致} 40 V／80 V·C。40 Vは合うがV·Cはエネルギーの単位で、電場ではない。\par

\noindent\textbf{4：条件に一致} 40 V／400 V/m。電位差の大きさ40 V、電場の大きさ400 V/mで両単位も合う。\par

\noindent\textbf{5：条件に不一致} 8 V／16 V·C。電位差は40 Vであり8 Vではなく、電場の単位も違う。\par

\paragraph{再利用する考え方} 電位・外力仕事と負電荷の符号を、成立条件・単位・符号から再現する。電場と電束密度、磁場と磁束密度、電荷一定と電圧一定を読み分ける。

\paragraph{共通理論} \hyperref[theory:em:potential]{第\ref*{theory:em:potential}節}

\par 参考：\cite{mit_em03}。

\paragraph{関連問題} \hyperref[q:r01-ret-applied-08]{R01-RET-APPLIED-08}

\par 難易度 2（校準初稿）、優先度未評価。

\subsection{平成28年度 問4}
\label{q:h28-reg-applied-04}

\noindent\textbf{原典：} sources/H28/17\_applied\_science.pdf、PDF物理2頁。

\textbf{公式正答：} 2（single）。

\paragraph{問題の要旨} 誘電体とコンデンサーについて、最も不適切な説明を選ぶ。

\Needspace{4\baselineskip}\paragraph{確認した条件・記述（独自要約）}\mbox{}\par

\noindent\textbf{方向：} 最も不適切。\par

\noindent\textbf{モデル：} 容量の比例関係は一様誘電体と平行平板の近似。\par

\paragraph{問われる核心} 誘電分極と強誘電性の区別

\paragraph{必要知識} 分極、容量、交流損失、電解コンデンサーの酸化膜。

\paragraph{考え方・正答を選ぶ理由（査読待ち初稿）} 電場による電気双極子の誘起・配向は分極であり、全ての誘電分極を強誘電性と呼ぶことはできない。外場0でも自発分極を持ち、それを電場で反転できることが強誘電体の特徴である。2を選ぶ。損失、材料、幾何学による容量変化は別の性質である。

\Needspace{4\baselineskip}\paragraph{選択肢の検討（査読待ち初稿）}\mbox{}\par

\noindent\textbf{1：適切} 静電容量の単位はファラドFである。。C=\allowbreak{}Q/Vなので単位はC/V=\allowbreak{}Fである。\par

\noindent\textbf{2：不適切} 絶縁体を電場へ置いて電気双極子モーメントを持つことを強誘電性という。。これは一般的な誘電分極である。強誘電性は外部電場なしの自発分極とその反転を特徴とする。\par

\noindent\textbf{3：適切} 交流電場で電気エネルギーの一部が熱となることを誘電損失という。。分極の応答に損失があると電場からのエネルギーが熱へ散逸する。\par

\noindent\textbf{4：適切} 可動電極の対向面積を変えると容量が変化する。。端効果を無視するとC=εA/dであり対向面積Aに依存する。\par

\noindent\textbf{5：適切} 誘電体として酸化アルミニウムや酸化タンタルが使われる。。金属酸化物の薄い絶縁層を使う電解コンデンサーの例である。\par

\paragraph{再利用する考え方} 誘電分極と強誘電性の区別を、成立条件・単位・符号から再現する。電場と電束密度、磁場と磁束密度、電荷一定と電圧一定を読み分ける。

\paragraph{共通理論} \hyperref[theory:em:dielectric]{第\ref*{theory:em:dielectric}節}, \hyperref[theory:em:capacitor]{第\ref*{theory:em:capacitor}節}

\par 参考：\cite{mit_em05,mit_ferroelectric}。

\paragraph{関連問題} \hyperref[q:r02-reg-applied-03]{R02-REG-APPLIED-03}

\par 難易度 2（校準初稿）、優先度未評価。

\subsection{平成29年度 問2}
\label{q:h29-reg-applied-02}

\noindent\textbf{原典：} sources/H29/17\_applied\_science.pdf、PDF物理1頁。

\textbf{公式正答：} 1（single）。

\paragraph{問題の要旨} 平行平板コンデンサーの容量と直列接続について、最も不適切なものを選ぶ。

\Needspace{4\baselineskip}\paragraph{確認した条件・記述（独自要約）}\mbox{}\par

\noindent\textbf{方向：} 最も不適切。\par

\noindent\textbf{モデル：} 平板の端効果を無視。容量比較では形状を固定。直列の2個は初め未充電。\par

\paragraph{問われる核心} 容量の幾何学と固定電荷

\paragraph{必要知識} 平行平板近似、比誘電率、V=\allowbreak{}Q/C、直列の電荷保存。

\paragraph{考え方・正答を選ぶ理由（査読待ち初稿）} \(C=\varepsilon A/d\)より誘電率、面積の増加と距離の減少は容量を大きくする。誘電率増加で容量が減るとする1が誤りである。Q一定の電圧はQ/C、未充電直列の同じ電気量は中間導体の電荷保存から確認する。

\Needspace{4\baselineskip}\paragraph{選択肢の検討（査読待ち初稿）}\mbox{}\par

\noindent\textbf{1：不適切} 比誘電率の大きな物質を極板間へ入れると容量は減少する。。C=εA/dなので同じ形状では誘電率の増加で容量は増え、減少しない。\par

\noindent\textbf{2：適切} 極板間距離を小さくすると容量は増える。。分母dが小さくなるとCが増える。\par

\noindent\textbf{3：適切} 極板面積に比例して容量は増える。。面積Aが増えるとCも同じ割合で増える。\par

\noindent\textbf{4：適切} 電気量が同じなら容量の小さい方が大きな電位差を持つ。。V=\allowbreak{}Q/CなのでQ一定ではCが小さいほどVが大きい。\par

\noindent\textbf{5：適切} 未充電の2個を直列につないで電圧を加えると等しい電気量が蓄えられる。。漏れのない直列回路で中間導体が初めに中性なら、電荷保存で対向極板の電荷の大きさは等しい。\par

\paragraph{再利用する考え方} 容量の幾何学と固定電荷を、成立条件・単位・符号から再現する。電場と電束密度、磁場と磁束密度、電荷一定と電圧一定を読み分ける。

\paragraph{共通理論} \hyperref[theory:em:capacitor]{第\ref*{theory:em:capacitor}節}, \hyperref[theory:em:dielectric]{第\ref*{theory:em:dielectric}節}

\par 参考：\cite{mit_em05}。

\paragraph{関連問題} \hyperref[q:r07-reg-applied-03]{R07-REG-APPLIED-03}

\par 難易度 2（校準初稿）、優先度未評価。

\subsection{令和元年度再試験 問8}
\label{q:r01-ret-applied-08}

\noindent\textbf{原典：} sources/R01\_retest/17\_applied\_science.pdf、PDF物理3頁。

\textbf{公式正答：} 3（single）。

\paragraph{問題の要旨} 負電荷を電場方向へゆっくり動かす外力の仕事から、電位差と電場の大きさを選ぶ。

\Needspace{4\baselineskip}\paragraph{確認した条件・記述（独自要約）}\mbox{}\par

\noindent\textbf{電荷：} −2 C\par

\noindent\textbf{移動：} PからQ。電場の方向と移動方向は同じ。\par

\noindent\textbf{外力の仕事：} 80 J。ゆっくり移動し運動エネルギー変化を無視。\par

\noindent\textbf{距離：} 0.2 m\par

\noindent\textbf{候補の順：} 電位差／電場。電位差は候補に合わせ大きさを表す。\par

\paragraph{問われる核心} 電位・外力仕事と負電荷の符号

\paragraph{必要知識} 準静的移動、位置エネルギーqV、E=−grad V、V/m=\allowbreak{}N/C。

\paragraph{考え方・正答を選ぶ理由（査読待ち初稿）} 外力の仕事は\(W_{ext}=q(V_Q-V_P)\)で、80=\allowbreak{}(−2)ΔVよりΔV=−40 Vである。電場方向へ進むと電位が下がるが、負電荷の位置エネルギーは増える。候補の電位差の大きさは40 V、\(|E|=|\Delta V|/d\)=200 V/mなので3である。電場の仕事は−80 Jと逆符号になる。

\Needspace{4\baselineskip}\paragraph{選択肢の検討（査読待ち初稿）}\mbox{}\par

\noindent\textbf{1：条件に不一致} 160 V／32 V·C。外力80 Jを電荷の絶対値2 Cで割ると40 Vであり、160 Vではない。\par

\noindent\textbf{2：条件に不一致} 160 V／800 V/m。電位差の大きさは40 Vで、電場も200 V/mである。\par

\noindent\textbf{3：条件に一致} 40 V／200 V/m。電位差の大きさ40 V、電場の大きさ200 V/mで両単位も合う。\par

\noindent\textbf{4：条件に不一致} 40 V／400 V/m。40 Vは合うが距離は0.2 mなので電場は200 V/mである。\par

\noindent\textbf{5：条件に不一致} 40 V／8 V·C。V·CはJに等しく、電場の単位V/mではない。\par

\paragraph{再利用する考え方} 電位・外力仕事と負電荷の符号を、成立条件・単位・符号から再現する。電場と電束密度、磁場と磁束密度、電荷一定と電圧一定を読み分ける。

\paragraph{共通理論} \hyperref[theory:em:potential]{第\ref*{theory:em:potential}節}

\par 参考：\cite{mit_em03}。

\paragraph{関連問題} \hyperref[q:h26-reg-applied-04]{H26-REG-APPLIED-04}

\par 難易度 2（校準初稿）、優先度未評価。

\subsection{令和2年度 問3}
\label{q:r02-reg-applied-03}

\noindent\textbf{原典：} sources/R02/17\_applied\_science.pdf、PDF物理2頁。

\textbf{公式正答：} 1（single）。

\paragraph{問題の要旨} 誘電体とコンデンサーについて、最も不適切な説明を選ぶ。

\Needspace{4\baselineskip}\paragraph{確認した条件・記述（独自要約）}\mbox{}\par

\noindent\textbf{方向：} 最も不適切。\par

\noindent\textbf{モデル：} 容量の比例関係は一様誘電体と平行平板の近似。\par

\paragraph{問われる核心} 誘電分極と強誘電性の区別

\paragraph{必要知識} 分極、容量、交流損失、電解コンデンサーの酸化膜。

\paragraph{考え方・正答を選ぶ理由（査読待ち初稿）} 電場による電気双極子の誘起・配向は分極であり、全ての誘電分極を強誘電性と呼ぶことはできない。外場0でも自発分極を持ち、それを電場で反転できることが強誘電体の特徴である。1を選ぶ。損失、材料、幾何学による容量変化は別の性質である。

\Needspace{4\baselineskip}\paragraph{選択肢の検討（査読待ち初稿）}\mbox{}\par

\noindent\textbf{1：不適切} 絶縁体を電場へ置いて電気双極子モーメントを持つことを強誘電性という。。これは一般的な誘電分極である。強誘電性は外部電場なしの自発分極とその反転を特徴とする。\par

\noindent\textbf{2：適切} 交流電場で電気エネルギーの一部が熱となることを誘電損失という。。分極の応答に損失があると電場からのエネルギーが熱へ散逸する。\par

\noindent\textbf{3：適切} 可動電極の対向面積を変えると容量が変化する。。端効果を無視するとC=εA/dであり対向面積Aに依存する。\par

\noindent\textbf{4：適切} 誘電体として酸化アルミニウムや酸化タンタルが使われる。。金属酸化物の薄い絶縁層を使う電解コンデンサーの例である。\par

\noindent\textbf{5：適切} 平行平板容量は誘電率に比例し、誘電体の厚さへ反比例する。。誘電体が一様に隙間を満たす平行平板近似ではC=εA/dである。\par

\paragraph{再利用する考え方} 誘電分極と強誘電性の区別を、成立条件・単位・符号から再現する。電場と電束密度、磁場と磁束密度、電荷一定と電圧一定を読み分ける。

\paragraph{共通理論} \hyperref[theory:em:dielectric]{第\ref*{theory:em:dielectric}節}, \hyperref[theory:em:capacitor]{第\ref*{theory:em:capacitor}節}

\par 参考：\cite{mit_em05,mit_ferroelectric}。

\paragraph{関連問題} \hyperref[q:h28-reg-applied-04]{H28-REG-APPLIED-04}

\par 難易度 2（校準初稿）、優先度未評価。

\subsection{令和5年度 問3}
\label{q:r05-reg-applied-03}

\noindent\textbf{原典：} sources/R05/17\_applied\_science.pdf、PDF物理2頁。

\textbf{公式正答：} 2（single）。

\paragraph{問題の要旨} 電荷を固定した平行平板の間隔と誘電率を変え、極板間引力の比を選ぶ。

\Needspace{4\baselineskip}\paragraph{確認した条件・記述（独自要約）}\mbox{}\par

\noindent\textbf{初期：} 真空、間隔d、帯電量Q、引力F0。\par

\noindent\textbf{変更：} 間隔5d、比誘電率5の誘電体で全間隙を満たす。\par

\noindent\textbf{固定量：} Qは変わらない。\par

\noindent\textbf{モデル：} 同じ面積A、端効果なし、一様な線形誘電体。\par

\paragraph{問われる核心} 固定電荷のエネルギー微分と静電力

\paragraph{必要知識} 容量、固定する量、平板近似、力とエネルギーの微分。

\paragraph{考え方・正答を選ぶ理由（査読待ち初稿）} 固定Qのエネルギーは\(U=Q^2/(2C)=Q^2d/(2\varepsilon A)\)であり、間隔方向の引力の大きさは\(F=\partial U/\partial d=Q^2/(2\varepsilon A)\)となる。εが5ε0へ変わるのでF=\allowbreak{}F0/5、2である。間隔5倍で容量は初期と同じになるが、力はエネルギーの間隔微分であり同じとは限らない。

\Needspace{4\baselineskip}\paragraph{選択肢の検討（査読待ち初稿）}\mbox{}\par

\noindent\textbf{1：条件に不一致} \(F=F_0/25\)。固定Qで引力は間隔に依存しないため、距離の逆二乗だけを使うのは違う。\par

\noindent\textbf{2：条件に一致} \(F=F_0/5\)。F=\allowbreak{}Q\textsuperscript{2}/(2εA)でεが5倍となり1/5である。\par

\noindent\textbf{3：条件に不一致} \(F=F_0\)。εの増加による引力減少を無視している。\par

\noindent\textbf{4：条件に不一致} \(F=5F_0\)。固定Qでは引力はεに反比例し、比例ではない。\par

\noindent\textbf{5：条件に不一致} \(F=25F_0\)。電荷固定の平板モデルでは25倍となる要因はない。\par

\paragraph{再利用する考え方} 固定電荷のエネルギー微分と静電力を、成立条件・単位・符号から再現する。電場と電束密度、磁場と磁束密度、電荷一定と電圧一定を読み分ける。

\paragraph{共通理論} \hyperref[theory:em:force]{第\ref*{theory:em:force}節}, \hyperref[theory:em:capacitor]{第\ref*{theory:em:capacitor}節}

\par 参考：\cite{mit_em05}。

\par 難易度 2（校準初稿）、優先度未評価。

\subsection{令和7年度 問3}
\label{q:r07-reg-applied-03}

\noindent\textbf{原典：} sources/R07/17\_applied\_science.pdf、PDF物理1頁。

\textbf{公式正答：} 3（single）。

\paragraph{問題の要旨} 平行平板コンデンサーの容量と直列接続について、最も不適切なものを選ぶ。

\Needspace{4\baselineskip}\paragraph{確認した条件・記述（独自要約）}\mbox{}\par

\noindent\textbf{方向：} 最も不適切。\par

\noindent\textbf{モデル：} 平板の端効果を無視。容量比較では形状を固定。直列の2個は初め未充電。\par

\paragraph{問われる核心} 容量の幾何学と固定電荷

\paragraph{必要知識} 平行平板近似、比誘電率、V=\allowbreak{}Q/C、直列の電荷保存。

\paragraph{考え方・正答を選ぶ理由（査読待ち初稿）} \(C=\varepsilon A/d\)より誘電率、面積の増加と距離の減少は容量を大きくする。誘電率増加で容量が減るとする3が誤りである。Q一定の電圧はQ/C、未充電直列の同じ電気量は中間導体の電荷保存から確認する。

\Needspace{4\baselineskip}\paragraph{選択肢の検討（査読待ち初稿）}\mbox{}\par

\noindent\textbf{1：適切} 極板面積に比例して容量は増える。。面積Aが増えるとCも同じ割合で増える。\par

\noindent\textbf{2：適切} 極板間距離を小さくすると容量は増える。。分母dが小さくなるとCが増える。\par

\noindent\textbf{3：不適切} 比誘電率の大きな物質を極板間へ入れると容量は減少する。。C=εA/dなので同じ形状では誘電率の増加で容量は増え、減少しない。\par

\noindent\textbf{4：適切} 未充電の2個を直列につないで電圧を加えると等しい電気量が蓄えられる。。漏れのない直列回路で中間導体が初めに中性なら、電荷保存で対向極板の電荷の大きさは等しい。\par

\noindent\textbf{5：適切} 電気量が同じなら容量の小さい方が大きな電位差を持つ。。V=\allowbreak{}Q/CなのでQ一定ではCが小さいほどVが大きい。\par

\paragraph{再利用する考え方} 容量の幾何学と固定電荷を、成立条件・単位・符号から再現する。電場と電束密度、磁場と磁束密度、電荷一定と電圧一定を読み分ける。

\paragraph{共通理論} \hyperref[theory:em:capacitor]{第\ref*{theory:em:capacitor}節}, \hyperref[theory:em:dielectric]{第\ref*{theory:em:dielectric}節}

\par 参考：\cite{mit_em05}。

\paragraph{関連問題} \hyperref[q:h29-reg-applied-02]{H29-REG-APPLIED-02}

\par 難易度 2（校準初稿）、優先度未評価。

\section{電磁誘導・磁性}

\subsection{平成27年度 問2}
\label{q:h27-reg-applied-02}

\noindent\textbf{原典：} sources/H27/17\_applied\_science.pdf、PDF物理1頁。

\textbf{公式正答：} 3（single）。

\paragraph{問題の要旨} コイルを貫く磁束が一定速度で変わるときの誘導起電力の大きさを求める。

\Needspace{4\baselineskip}\paragraph{確認した条件・記述（独自要約）}\mbox{}\par

\noindent\textbf{巻数：} 2000\par

\noindent\textbf{磁束変化：} 各巻きを貫く磁束が20秒で0.01 Wb変化。\par

\noindent\textbf{速度：} 一定の割合。1 Wb=\allowbreak{}1 V·s。\par

\paragraph{問われる核心} 磁束鎖交数の変化と誘導

\paragraph{必要知識} ファラデーの法則、巻数、Wb/s=\allowbreak{}V。

\paragraph{考え方・正答を選ぶ理由（査読待ち初稿）} \(|\mathcal E|=N|\Delta\Phi|/\Delta t=2000(0.01)/20=1.0\) Vで3を選ぶ。一定速度なので平均変化率が任意の時点の変化率に等しい。向きは磁束変化を妨げるレンツの法則で決まり、ここでは大きさのみを求める。

\Needspace{4\baselineskip}\paragraph{選択肢の検討（査読待ち初稿）}\mbox{}\par

\noindent\textbf{1：条件に不一致} 0.10 V。2000×0.01/20=\allowbreak{}1 Vと違う。\par

\noindent\textbf{2：条件に不一致} 0.50 V。正しい値の1/2であり余分な係数がある。\par

\noindent\textbf{3：条件に一致} 1.0 V。巻数×各巻きの磁束変化率に一致する。\par

\noindent\textbf{4：条件に不一致} 2.5 V。磁束変化の時間20 sで割った1 Vと違う。\par

\noindent\textbf{5：条件に不一致} 5.0 V。与えられた変化率では1 Vである。\par

\paragraph{再利用する考え方} 磁束鎖交数の変化と誘導を、成立条件・単位・符号から再現する。電場と電束密度、磁場と磁束密度、電荷一定と電圧一定を読み分ける。

\paragraph{共通理論} \hyperref[theory:em:induction]{第\ref*{theory:em:induction}節}

\par 参考：\cite{mit_em10}。

\par 難易度 2（校準初稿）、優先度未評価。

\subsection{平成28年度 問3}
\label{q:h28-reg-applied-03}

\noindent\textbf{原典：} sources/H28/17\_applied\_science.pdf、PDF物理1頁。

\textbf{公式正答：} 3（single）。

\paragraph{問題の要旨} 電束密度・磁束密度を用いた電磁気の法則で、最も不適切なものを選ぶ。

\Needspace{4\baselineskip}\paragraph{確認した条件・記述（独自要約）}\mbox{}\par

\noindent\textbf{記号：} E=電場、D=電束密度、H=磁場、B=磁束密度、ρ=真電荷密度、i=伝導電流密度。\par

\noindent\textbf{方向：} 最も不適切な記述。\par

\paragraph{問われる核心} マクスウェル方程式と場の発散

\paragraph{必要知識} 自由電荷と束縛電荷、回転と発散、変位電流、磁気単極子。

\paragraph{考え方・正答を選ぶ理由（査読待ち初稿）} 磁場に関するガウス法則\(\operatorname{div}B=0\)は、閉曲面を通る正味磁束が0であること、すなわち通常の理論で磁気単極子がないことを表す。「存在」を反映するとする3が誤りである。他の式は真電荷、磁束の変化、変位電流をそれぞれ扱う。

\Needspace{4\baselineskip}\paragraph{選択肢の検討（査読待ち初稿）}\mbox{}\par

\noindent\textbf{1：適切} 閉曲面を貫く電束は、内部にある真電荷の総量に等しい。。Dの閉曲面電束は自由電荷、原典の真電荷の総量に等しい。束縛電荷まで含むEの式と区別する。\par

\noindent\textbf{2：適切} \(\operatorname{rot}E=-\partial B/\partial t\)は時間変化する磁束密度が電場を誘導する法則である。。ファラデーの法則の局所形で、負号は向きを表す。\par

\noindent\textbf{3：不適切} \(\operatorname{div}B=0\)は磁気単極子の存在を反映する。。発散0は磁気単極子による湧き出し・吸込みがないという通常の電磁気学の式であり、存在と逆である。\par

\noindent\textbf{4：適切} \(\operatorname{rot}H=i+\partial D/\partial t\)で\(\partial D/\partial t\)は変位電流密度である。。伝導電流に加える変位電流の項を正しく述べる。\par

\noindent\textbf{5：適切} 磁場のガウス法則は時間変化する磁場にも成立する。。div B=\allowbreak{}0は静磁場だけの制限ではなく、マクスウェル方程式の一つである。\par

\paragraph{再利用する考え方} マクスウェル方程式と場の発散を、成立条件・単位・符号から再現する。電場と電束密度、磁場と磁束密度、電荷一定と電圧一定を読み分ける。

\paragraph{共通理論} \hyperref[theory:em:maxwell]{第\ref*{theory:em:maxwell}節}, \hyperref[theory:em:gauss]{第\ref*{theory:em:gauss}節}

\par 参考：\cite{mit_em13}。

\paragraph{関連問題} \hyperref[q:r03-reg-applied-04]{R03-REG-APPLIED-04}

\par 難易度 2（校準初稿）、優先度未評価。

\subsection{平成30年度 問4}
\label{q:h30-reg-applied-04}

\noindent\textbf{原典：} sources/H30/17\_applied\_science.pdf、PDF物理2頁。

\textbf{公式正答：} 3（single）。

\paragraph{問題の要旨} 小さな電流素片が作る磁場を与える法則名を選ぶ。

\Needspace{4\baselineskip}\paragraph{確認した条件・記述（独自要約）}\mbox{}\par

\noindent\textbf{対象：} 定常電流の微小素片が作る磁場。\par

\paragraph{問われる核心} 電流素片から磁場を作る法則

\paragraph{必要知識} 定常電流、外積、重ね合わせ、磁束の時間変化との違い。

\paragraph{考え方・正答を選ぶ理由（査読待ち初稿）} 微小素片の寄与は\(dB=\mu_0 I(d\ell\times\hat r)/(4\pi r^2)\)というビオ・サバールの法則で与えられ、3である。素片を積分して導線全体の場を求める。誘導起電力の法則と、電流を源とする場の法則を区別する。

\Needspace{4\baselineskip}\paragraph{選択肢の検討（査読待ち初稿）}\mbox{}\par

\noindent\textbf{1：条件に不一致} クーロンの法則。点電荷の静電気力を与え、電流素片の磁場ではない。\par

\noindent\textbf{2：条件に不一致} プランクの法則。黒体の熱放射スペクトルの法則である。\par

\noindent\textbf{3：条件に一致} ビオ・サバールの法則。定常電流素片Idℓから磁束密度dBを与える法則である。\par

\noindent\textbf{4：条件に不一致} ファラデーの誘導法則。磁束の時間変化による誘導起電力を与える。\par

\noindent\textbf{5：条件に不一致} キュリー・ワイスの法則。磁性体の磁化率の温度依存を扱う。\par

\paragraph{再利用する考え方} 電流素片から磁場を作る法則を、成立条件・単位・符号から再現する。電場と電束密度、磁場と磁束密度、電荷一定と電圧一定を読み分ける。

\paragraph{共通理論} \hyperref[theory:em:biot]{第\ref*{theory:em:biot}節}

\par 参考：\cite{mit_em09}。

\paragraph{関連問題} \hyperref[q:r04-reg-applied-03]{R04-REG-APPLIED-03}

\par 難易度 2（校準初稿）、優先度未評価。

\subsection{令和元年度 問4}
\label{q:r01-reg-applied-04}

\noindent\textbf{原典：} sources/R01/17\_applied\_science.pdf、PDF物理2頁。

\textbf{公式正答：} 3（single）。

\paragraph{問題の要旨} 電磁気の基本法則の説明で、最も不適切なものを選ぶ。

\Needspace{4\baselineskip}\paragraph{確認した条件・記述（独自要約）}\mbox{}\par

\noindent\textbf{注意：} ③の原典は二箇所とも電束と記す。磁束へ書き換えて読まない。\par

\noindent\textbf{方向：} 最も不適切。\par

\paragraph{問われる核心} 電束の変化と磁束の変化

\paragraph{必要知識} ガウス、ビオ・サバール、ファラデー、連続の式、変位電流。

\paragraph{考え方・正答を選ぶ理由（査読待ち初稿）} 3の決め手は原典の「電束」である。ファラデーの法則は\(\mathcal E=-d\Phi_B/dt\)と磁束を用いる。電束変化による変位電流と混同しているため3を選ぶ。保存則の「出入り」は流入を正にした正味量として区別する。

\Needspace{4\baselineskip}\paragraph{選択肢の検討（査読待ち初稿）}\mbox{}\par

\noindent\textbf{1：適切} 閉曲面を貫く電束は内部の真電荷の量に等しく、電場のガウス法則である。。Dの閉曲面電束を内部真電荷へ対応させる説明である。\par

\noindent\textbf{2：適切} ビオ・サバールの法則は定常電流が作る磁場を決める。。電流素片から場を積分する定常磁場の法則である。\par

\noindent\textbf{3：不適切} 閉回路を貫く電束が変化すると、電束変化を妨げる起電力が生じることをファラデーの法則という。。誘導起電力は磁束の変化に対応する。電束変化はアンペール・マクスウェルの変位電流側の量である。\par

\noindent\textbf{4：適切} 領域の総電荷の時間変化は閉曲面を通じた電荷の正味の出入りと対応し、電気量保存を表す。。外向き電流を正に取ればdQ/dt=−∫i·n dSである。原典の出入りは符号付きの正味量として読む。\par

\noindent\textbf{5：適切} マクスウェル方程式は電磁場の基本方程式で19世紀に提出された。。マクスウェルが19世紀に電磁気の法則をまとめたという説明に合う。\par

\paragraph{再利用する考え方} 電束の変化と磁束の変化を、成立条件・単位・符号から再現する。電場と電束密度、磁場と磁束密度、電荷一定と電圧一定を読み分ける。

\paragraph{共通理論} \hyperref[theory:em:induction]{第\ref*{theory:em:induction}節}, \hyperref[theory:em:maxwell]{第\ref*{theory:em:maxwell}節}

\par 参考：\cite{mit_em10,mit_em13}。

\par 難易度 2（校準初稿）、優先度未評価。

\subsection{令和3年度 問4}
\label{q:r03-reg-applied-04}

\noindent\textbf{原典：} sources/R03/17\_applied\_science.pdf、PDF物理2頁。

\textbf{公式正答：} 2（single）。

\paragraph{問題の要旨} 電束密度・磁束密度を用いた電磁気の法則で、最も不適切なものを選ぶ。

\Needspace{4\baselineskip}\paragraph{確認した条件・記述（独自要約）}\mbox{}\par

\noindent\textbf{記号：} E=電場、D=電束密度、H=磁場、B=磁束密度、ρ=真電荷密度、i=伝導電流密度。\par

\noindent\textbf{方向：} 最も不適切な記述。\par

\paragraph{問われる核心} マクスウェル方程式と場の発散

\paragraph{必要知識} 自由電荷と束縛電荷、回転と発散、変位電流、磁気単極子。

\paragraph{考え方・正答を選ぶ理由（査読待ち初稿）} 磁場に関するガウス法則\(\operatorname{div}B=0\)は、閉曲面を通る正味磁束が0であること、すなわち通常の理論で磁気単極子がないことを表す。「存在」を反映するとする2が誤りである。他の式は真電荷、磁束の変化、変位電流をそれぞれ扱う。

\Needspace{4\baselineskip}\paragraph{選択肢の検討（査読待ち初稿）}\mbox{}\par

\noindent\textbf{1：適切} \(\operatorname{rot}E=-\partial B/\partial t\)は時間変化する磁束密度が電場を誘導する法則である。。ファラデーの法則の局所形で、負号は向きを表す。\par

\noindent\textbf{2：不適切} \(\operatorname{div}B=0\)は磁気単極子の存在を反映する。。発散0は磁気単極子による湧き出し・吸込みがないという通常の電磁気学の式であり、存在と逆である。\par

\noindent\textbf{3：適切} 磁場のガウス法則は時間変化する磁場にも成立する。。div B=\allowbreak{}0は静磁場だけの制限ではなく、マクスウェル方程式の一つである。\par

\noindent\textbf{4：適切} 閉曲面を貫く電束は、内部にある真電荷の総量に等しい。。Dの閉曲面電束は自由電荷、原典の真電荷の総量に等しい。束縛電荷まで含むEの式と区別する。\par

\noindent\textbf{5：適切} \(\operatorname{rot}H=i+\partial D/\partial t\)で\(\partial D/\partial t\)は変位電流密度である。。伝導電流に加える変位電流の項を正しく述べる。\par

\paragraph{再利用する考え方} マクスウェル方程式と場の発散を、成立条件・単位・符号から再現する。電場と電束密度、磁場と磁束密度、電荷一定と電圧一定を読み分ける。

\paragraph{共通理論} \hyperref[theory:em:maxwell]{第\ref*{theory:em:maxwell}節}, \hyperref[theory:em:gauss]{第\ref*{theory:em:gauss}節}

\par 参考：\cite{mit_em13}。

\paragraph{関連問題} \hyperref[q:h28-reg-applied-03]{H28-REG-APPLIED-03}

\par 難易度 2（校準初稿）、優先度未評価。

\subsection{令和4年度 問3}
\label{q:r04-reg-applied-03}

\noindent\textbf{原典：} sources/R04/17\_applied\_science.pdf、PDF物理1頁。

\textbf{公式正答：} 4（single）。

\paragraph{問題の要旨} 小さな電流素片が作る磁場を与える法則名を選ぶ。

\Needspace{4\baselineskip}\paragraph{確認した条件・記述（独自要約）}\mbox{}\par

\noindent\textbf{対象：} 定常電流の微小素片が作る磁場。\par

\paragraph{問われる核心} 電流素片から磁場を作る法則

\paragraph{必要知識} 定常電流、外積、重ね合わせ、磁束の時間変化との違い。

\paragraph{考え方・正答を選ぶ理由（査読待ち初稿）} 微小素片の寄与は\(dB=\mu_0 I(d\ell\times\hat r)/(4\pi r^2)\)というビオ・サバールの法則で与えられ、4である。素片を積分して導線全体の場を求める。誘導起電力の法則と、電流を源とする場の法則を区別する。

\Needspace{4\baselineskip}\paragraph{選択肢の検討（査読待ち初稿）}\mbox{}\par

\noindent\textbf{1：条件に不一致} プランクの法則。黒体の熱放射スペクトルの法則である。\par

\noindent\textbf{2：条件に不一致} クーロンの法則。点電荷の静電気力を与え、電流素片の磁場ではない。\par

\noindent\textbf{3：条件に不一致} ファラデーの誘導法則。磁束の時間変化による誘導起電力を与える。\par

\noindent\textbf{4：条件に一致} ビオ・サバールの法則。定常電流素片Idℓから磁束密度dBを与える法則である。\par

\noindent\textbf{5：条件に不一致} キュリー・ワイスの法則。磁性体の磁化率の温度依存を扱う。\par

\paragraph{再利用する考え方} 電流素片から磁場を作る法則を、成立条件・単位・符号から再現する。電場と電束密度、磁場と磁束密度、電荷一定と電圧一定を読み分ける。

\paragraph{共通理論} \hyperref[theory:em:biot]{第\ref*{theory:em:biot}節}

\par 参考：\cite{mit_em09}。

\paragraph{関連問題} \hyperref[q:h30-reg-applied-04]{H30-REG-APPLIED-04}

\par 難易度 2（校準初稿）、優先度未評価。

\subsection{令和6年度 問4}
\label{q:r06-reg-applied-04}

\noindent\textbf{原典：} sources/R06/17\_applied\_science.pdf、PDF物理2頁。

\textbf{公式正答：} 2（single）。

\paragraph{問題の要旨} 静電場、磁場中の運動、誘導、エネルギー流の説明で最も不適切なものを選ぶ。

\Needspace{4\baselineskip}\paragraph{確認した条件・記述（独自要約）}\mbox{}\par

\noindent\textbf{方向：} 最も不適切。\par

\noindent\textbf{導体：} 孤立導体球を帯電させた静電平衡。\par

\noindent\textbf{運動：} 磁場に垂直な初速度、磁場以外の力なし。\par

\paragraph{問われる核心} 静電平衡導体の内外と電磁場

\paragraph{必要知識} 導体遮蔽、クーロン則、ローレンツ力、磁束、外積と流密度。

\paragraph{考え方・正答を選ぶ理由（査読待ち初稿）} 導体内に非零の電場が残ると自由電荷が移動して平衡ではなくなる。静電平衡では内部電場は0で、表面の外側には点電荷に相当する逆二乗場がある。内外を入れ替えた2を選ぶ。磁気力は仕事をせず、エネルギー輸送はポインティングベクトルで記述する。

\Needspace{4\baselineskip}\paragraph{選択肢の検討（査読待ち初稿）}\mbox{}\par

\noindent\textbf{1：適切} 点電荷間の力は距離の二乗に反比例する。。クーロン力の逆二乗則に一致する。\par

\noindent\textbf{2：不適切} 帯電導体球の内部電場は中心からの距離の二乗に反比例する。。静電平衡の導体内部はE=\allowbreak{}0である。逆二乗場は孤立球の外側の説明である。\par

\noindent\textbf{3：適切} 磁場に垂直な速度を持つ荷電粒子は、速度に垂直な力で円運動する。。qv×Bは速度へ垂直で速さを変えず、向きを変える。\par

\noindent\textbf{4：適切} 起電力の大きさは回路を貫く磁束の時間変化率に比例する。。ファラデーの法則の大きさに一致する。\par

\noindent\textbf{5：適切} ポインティングベクトルはエネルギー流の密度で、電場・磁場に垂直である。。S=\allowbreak{}E×Hでエネルギー流密度W/m\textsuperscript{2}を表し、外積なので双方へ垂直である。\par

\paragraph{再利用する考え方} 静電平衡導体の内外と電磁場を、成立条件・単位・符号から再現する。電場と電束密度、磁場と磁束密度、電荷一定と電圧一定を読み分ける。

\paragraph{共通理論} \hyperref[theory:em:gauss]{第\ref*{theory:em:gauss}節}, \hyperref[theory:em:lorentz]{第\ref*{theory:em:lorentz}節}, \hyperref[theory:em:poynting]{第\ref*{theory:em:poynting}節}

\par 参考：\cite{mit_em04,mit_em13,mit_lorentz}。

\par 難易度 2（校準初稿）、優先度未評価。

\subsection{令和7年度 問4}
\label{q:r07-reg-applied-04}

\noindent\textbf{原典：} sources/R07/17\_applied\_science.pdf、PDF物理2頁。

\textbf{公式正答：} 4（single）。

\paragraph{問題の要旨} 磁性体で満たしたソレノイド内部の磁束密度を計算する。

\Needspace{4\baselineskip}\paragraph{確認した条件・記述（独自要約）}\mbox{}\par

\noindent\textbf{長さ：} 40 cm=\allowbreak{}0.40 m\par

\noindent\textbf{全巻数：} 150\par

\noindent\textbf{比透磁率：} 100\par

\noindent\textbf{電流：} 1.5 A\par

\noindent\textbf{磁気定数：} \(4\pi\times10^{-7}\) T·m/A\par

\noindent\textbf{モデル：} 長いソレノイドの内部近似。比透磁率は与えられた値を使用。\par

\paragraph{問われる核心} ソレノイドの巻数密度と磁束密度

\paragraph{必要知識} アンペール則、透磁率、全巻数と巻数密度、cmとm、TとmT。

\paragraph{考え方・正答を選ぶ理由（査読待ち初稿）} 単位長さの巻数は150/0.40=\allowbreak{}375 m\textsuperscript{-1}である。\(B=\mu_0\mu_r(N/\ell)I=4\pi\times10^{-7}\times100\times375\times1.5\)≃0.0707 T=\allowbreak{}70.7 mTなので4を選ぶ。端部や磁性体の非線形性はこの出題の近似へ含めず、与えられた比透磁率を用いる。

\Needspace{4\baselineskip}\paragraph{選択肢の検討（査読待ち初稿）}\mbox{}\par

\noindent\textbf{1：条件に不一致} \(2.8\times10^{-3}\) mT。与えられた透磁率・巻数密度では約70.7 mTであり一致しない。\par

\noindent\textbf{2：条件に不一致} \(1.4\times10^{-1}\) mT。透磁率と巻数密度を含めた70.7 mTと一致しない。\par

\noindent\textbf{3：条件に不一致} \(1.1\times10^1\) mT。約70.7 mTより小さく、全巻数を単位長さ巻数へ換算する必要がある。\par

\noindent\textbf{4：条件に一致} \(7.1\times10^1\) mT。B=μ0μr(N/ℓ)Iから約70.7 mTで最も近い。\par

\noindent\textbf{5：条件に不一致} \(1.1\times10^2\) mT。同じ計算では約70.7 mTで110 mTではない。\par

\paragraph{再利用する考え方} ソレノイドの巻数密度と磁束密度を、成立条件・単位・符号から再現する。電場と電束密度、磁場と磁束密度、電荷一定と電圧一定を読み分ける。

\paragraph{共通理論} \hyperref[theory:em:solenoid]{第\ref*{theory:em:solenoid}節}

\par 参考：\cite{mit_em09}。

\par 難易度 2（校準初稿）、優先度未評価。

\section{回路・抵抗}

\subsection{平成24年度 問4}
\label{q:h24-reg-applied-04}

\noindent\textbf{原典：} sources/H24/17\_applied\_science.pdf、PDF物理1頁。

\textbf{公式正答：} 3（single）。

\paragraph{問題の要旨} 10 Ωと5 Ωを並列接続した合成抵抗を選ぶ。

\Needspace{4\baselineskip}\paragraph{確認した条件・記述（独自要約）}\mbox{}\par

\noindent\textbf{接続：} 同じ二端子間へ並列に接続。\par

\paragraph{問われる核心} 並列抵抗の電流加算

\paragraph{必要知識} オーム則、コンダクタンス、直列と並列。

\paragraph{考え方・正答を選ぶ理由（査読待ち初稿）} \(1/R=1/10+1/5=3/10\)よりR=\allowbreak{}10/3 Ω≃3.3 Ωで3を選ぶ。並列では電圧が共通、電流が加算されるため、抵抗そのものではなくコンダクタンスを足す。

\Needspace{4\baselineskip}\paragraph{選択肢の検討（査読待ち初稿）}\mbox{}\par

\noindent\textbf{1：条件に不一致} 約0.2 Ω。0.2は5 Ωの逆数だが単位はSで、抵抗ではない。\par

\noindent\textbf{2：条件に不一致} 約1.5 Ω。逆数の和を逆に戻した10/3 Ωと違う。\par

\noindent\textbf{3：条件に一致} 約3.3 Ω。共通電圧で各電流を足した等価抵抗10/3 Ωに最も近い。\par

\noindent\textbf{4：条件に不一致} 約7.5 Ω。合成抵抗が小さい方の5 Ωを超えており並列の性質に反する。\par

\noindent\textbf{5：条件に不一致} 約15 Ω。10+5は直列接続の場合である。\par

\paragraph{再利用する考え方} 並列抵抗の電流加算を、成立条件・単位・符号から再現する。電場と電束密度、磁場と磁束密度、電荷一定と電圧一定を読み分ける。

\paragraph{共通理論} \hyperref[theory:em:dc]{第\ref*{theory:em:dc}節}

\par 参考：\cite{mit_em07}。

\paragraph{関連問題} \hyperref[q:h26-reg-applied-02]{H26-REG-APPLIED-02}

\par 難易度 2（校準初稿）、優先度未評価。

\subsection{平成26年度 問2}
\label{q:h26-reg-applied-02}

\noindent\textbf{原典：} sources/H26/17\_applied\_science.pdf、PDF物理1頁。

\textbf{公式正答：} 3（single）。

\paragraph{問題の要旨} 10 Ωと5 Ωを並列接続した合成抵抗を選ぶ。

\Needspace{4\baselineskip}\paragraph{確認した条件・記述（独自要約）}\mbox{}\par

\noindent\textbf{接続：} 同じ二端子間へ並列に接続。\par

\paragraph{問われる核心} 並列抵抗の電流加算

\paragraph{必要知識} オーム則、コンダクタンス、直列と並列。

\paragraph{考え方・正答を選ぶ理由（査読待ち初稿）} \(1/R=1/10+1/5=3/10\)よりR=\allowbreak{}10/3 Ω≃3.3 Ωで3を選ぶ。並列では電圧が共通、電流が加算されるため、抵抗そのものではなくコンダクタンスを足す。

\Needspace{4\baselineskip}\paragraph{選択肢の検討（査読待ち初稿）}\mbox{}\par

\noindent\textbf{1：条件に不一致} 約0.2 Ω。0.2は5 Ωの逆数だが単位はSで、抵抗ではない。\par

\noindent\textbf{2：条件に不一致} 約1.5 Ω。逆数の和を逆に戻した10/3 Ωと違う。\par

\noindent\textbf{3：条件に一致} 約3.3 Ω。共通電圧で各電流を足した等価抵抗10/3 Ωに最も近い。\par

\noindent\textbf{4：条件に不一致} 約7.5 Ω。合成抵抗が小さい方の5 Ωを超えており並列の性質に反する。\par

\noindent\textbf{5：条件に不一致} 約15 Ω。10+5は直列接続の場合である。\par

\paragraph{再利用する考え方} 並列抵抗の電流加算を、成立条件・単位・符号から再現する。電場と電束密度、磁場と磁束密度、電荷一定と電圧一定を読み分ける。

\paragraph{共通理論} \hyperref[theory:em:dc]{第\ref*{theory:em:dc}節}

\par 参考：\cite{mit_em07}。

\paragraph{関連問題} \hyperref[q:h24-reg-applied-04]{H24-REG-APPLIED-04}

\par 難易度 2（校準初稿）、優先度未評価。

\subsection{平成29年度 問3}
\label{q:h29-reg-applied-03}

\noindent\textbf{原典：} sources/H29/17\_applied\_science.pdf、PDF物理1頁。

\textbf{公式正答：} 3（single）。

\paragraph{問題の要旨} 直列RLC回路のインピーダンスの大きさが最小となる交流周波数を求める。

\Needspace{4\baselineskip}\paragraph{確認した条件・記述（独自要約）}\mbox{}\par

\noindent\textbf{抵抗：} 100 Ω\par

\noindent\textbf{インダクタンス：} 1 mH\par

\noindent\textbf{容量：} 0.1 μF\par

\noindent\textbf{接続：} 抵抗・コイル・コンデンサーを直列。理想的な周波数一定素子の定常交流。\par

\paragraph{問われる核心} 直列RLCのリアクタンス相殺

\paragraph{必要知識} 複素インピーダンス、角周波数と周波数、m=\allowbreak{}10\textsuperscript{-3}、μ=10\textsuperscript{-6}。

\paragraph{考え方・正答を選ぶ理由（査読待ち初稿）} \(Z=R+j(\omega L-1/(\omega C))\)なので大きさの二乗は\(R^2+(\omega L-1/(\omega C))^2\)となる。非負の第2項が0となる\(f_0=1/(2\pi\sqrt{LC})\)で最小である。mHとμFをHとFへ換算すると約15.9155 kHz、3を選ぶ。Rは最小値を決めるが、この理想モデルの共振周波数は変えない。

\Needspace{4\baselineskip}\paragraph{選択肢の検討（査読待ち初稿）}\mbox{}\par

\noindent\textbf{1：条件に不一致} 1 kHz。SIへ換算した共振周波数は約15.9155 kHzであり、この候補より離れる。\par

\noindent\textbf{2：条件に不一致} 5.1 kHz。SIへ換算した共振周波数は約15.9155 kHzであり、この候補より離れる。\par

\noindent\textbf{3：条件に一致} 16 kHz。リアクタンスが相殺する約15.9155 kHzに最も近い。\par

\noindent\textbf{4：条件に不一致} 100 kHz。SIへ換算した共振周波数は約15.9155 kHzであり、この候補より離れる。\par

\noindent\textbf{5：条件に不一致} 630 kHz。SIへ換算した共振周波数は約15.9155 kHzであり、この候補より離れる。\par

\paragraph{再利用する考え方} 直列RLCのリアクタンス相殺を、成立条件・単位・符号から再現する。電場と電束密度、磁場と磁束密度、電荷一定と電圧一定を読み分ける。

\paragraph{共通理論} \hyperref[theory:em:rlc]{第\ref*{theory:em:rlc}節}

\par 参考：\cite{mit_em12}。

\paragraph{関連問題} \hyperref[q:r02-reg-applied-04]{R02-REG-APPLIED-04}

\par 難易度 2（校準初稿）、優先度未評価。

\subsection{平成30年度 問3}
\label{q:h30-reg-applied-03}

\noindent\textbf{原典：} sources/H30/17\_applied\_science.pdf、PDF物理2頁。

\textbf{公式正答：} 3（single）。

\paragraph{問題の要旨} 電流と電気抵抗について、最も不適切な記述を選ぶ。

\Needspace{4\baselineskip}\paragraph{確認した条件・記述（独自要約）}\mbox{}\par

\noindent\textbf{対象：} 一定の抵抗率を持つ一様な導体。オーム則を使う際は温度等を固定。\par

\paragraph{問われる核心} 抵抗率と導体の幾何学

\paragraph{必要知識} 電流の定義、電子の電荷、オーム則、ジュール電力。

\paragraph{考え方・正答を選ぶ理由（査読待ち初稿）} 一様導体では長いほど衝突する経路が増え、太いほど電流の経路が増える。\(R=\rho L/A\)なので3の比例関係が逆である。電力と電流の関係はオーム性・温度一定という条件で読む。

\Needspace{4\baselineskip}\paragraph{選択肢の検討（査読待ち初稿）}\mbox{}\par

\noindent\textbf{1：適切} 電流は断面を単位時間に通過する電気量である。。I=\allowbreak{}dQ/dtという定義に一致する。\par

\noindent\textbf{2：適切} 自由電子は負電荷なので電場と逆向きの力を受けて流れる。。電気力qEでq<0なので力はEと逆向きである。\par

\noindent\textbf{3：不適切} 抵抗は断面積に比例し長さに反比例する。。R=ρL/Aであり、長さと断面積の比例・反比例が逆である。\par

\noindent\textbf{4：適切} 電力は単位時間の仕事で、電圧の二乗に比例し抵抗へ反比例する。。抵抗性負荷ではP=\allowbreak{}VI=\allowbreak{}V\textsuperscript{2}/Rとなる。\par

\noindent\textbf{5：適切} 電流は電圧に比例し、抵抗が大きいほど流れにくい。。オーム性導体のI=\allowbreak{}V/Rに一致する。\par

\paragraph{再利用する考え方} 抵抗率と導体の幾何学を、成立条件・単位・符号から再現する。電場と電束密度、磁場と磁束密度、電荷一定と電圧一定を読み分ける。

\paragraph{共通理論} \hyperref[theory:em:dc]{第\ref*{theory:em:dc}節}, \hyperref[theory:em:resistivity]{第\ref*{theory:em:resistivity}節}

\par 参考：\cite{mit_em06}。

\par 難易度 2（校準初稿）、優先度未評価。

\subsection{令和元年度 問3}
\label{q:r01-reg-applied-03}

\noindent\textbf{原典：} sources/R01/17\_applied\_science.pdf、PDF物理1頁。

\textbf{公式正答：} 5（single）。

\paragraph{問題の要旨} 電流の定義、向き、磁場、単位の4空欄の組合せを選ぶ。

\Needspace{4\baselineskip}\paragraph{確認した条件・記述（独自要約）}\mbox{}\par

\noindent\textbf{A：} 電流が表す何の流れか。\par

\noindent\textbf{B：} 電流方向と逆に移動するAの符号。\par

\noindent\textbf{C：} 電流方向に対する、直線電流の周囲の磁場方向。\par

\noindent\textbf{D：} アンペアの表現。\par

\noindent\textbf{並び：} A／B／C／D。\par

\paragraph{問われる核心} 電流の符号規約と単位

\paragraph{必要知識} 電荷輸送、正負の担い手、右ねじ、アンペア。

\paragraph{考え方・正答を選ぶ理由（査読待ち初稿）} 電流は電荷の移動を数える。正電荷の移動方向を正とするため、負電荷の移動と電流は逆向きになる。直線電流の周囲の磁場は電流を囲む円周方向、1 A=\allowbreak{}1 C/sで5を選ぶ。電流の担い手を電子だけへ限定しない。

\Needspace{4\baselineskip}\paragraph{選択肢の検討（査読待ち初稿）}\mbox{}\par

\noindent\textbf{1：条件に不一致} 原子／負／垂直／C/s。電流は一般に電荷の流れで、原子に限定されない。\par

\noindent\textbf{2：条件に不一致} 電子／正／垂直／C/s。電子は負であり、正の電子が逆方向へ進むという表現は合わない。\par

\noindent\textbf{3：条件に不一致} イオン／正／平行／C/m\textsuperscript{2}。正電荷の移動方向は電流方向と同じで、磁場と単位も違う。\par

\noindent\textbf{4：条件に不一致} 電荷／負／平行／C/m\textsuperscript{2}。最初の二語は合うが、直線電流の磁場は周方向で電流へ垂直、単位はC/sである。\par

\noindent\textbf{5：条件に一致} 電荷／負／垂直／C/s。電荷の流れ、負の移動と逆向き、垂直、C/sの4条件が全て合う。\par

\paragraph{再利用する考え方} 電流の符号規約と単位を、成立条件・単位・符号から再現する。電場と電束密度、磁場と磁束密度、電荷一定と電圧一定を読み分ける。

\paragraph{共通理論} \hyperref[theory:em:dc]{第\ref*{theory:em:dc}節}, \hyperref[theory:em:biot]{第\ref*{theory:em:biot}節}

\par 参考：\cite{mit_em06,mit_em09}。

\par 難易度 2（校準初稿）、優先度未評価。

\subsection{令和元年度再試験 問6}
\label{q:r01-ret-applied-06}

\noindent\textbf{原典：} sources/R01\_retest/17\_applied\_science.pdf、PDF物理3頁。

\textbf{公式正答：} 4（single）。

\paragraph{問題の要旨} 直流回路と抵抗の説明で、最も不適切なものを選ぶ。

\Needspace{4\baselineskip}\paragraph{確認した条件・記述（独自要約）}\mbox{}\par

\noindent\textbf{対象：} 定常直流、真性半導体、純金属の抵抗。\par

\noindent\textbf{方向：} 最も不適切。\par

\paragraph{問われる核心} 抵抗の温度依存と材料の区別

\paragraph{必要知識} 金属の散乱、残留抵抗、真性キャリア、近似の温度範囲。

\paragraph{考え方・正答を選ぶ理由（査読待ち初稿）} 4は材料と温度域を限定せず二乗比例を一般則とする点が誤りである。金属の室温近傍では\(\rho(T)\simeq\rho(T_0)[1+\alpha(T-T_0)]\)と近似し、真性半導体では担い手の熱励起が支配する。低温でT\textsuperscript{2}項が現れる場合まで否定するのではない。

\Needspace{4\baselineskip}\paragraph{選択肢の検討（査読待ち初稿）}\mbox{}\par

\noindent\textbf{1：適切} 二点間の電位差はどの経路で測っても同じである。。静電的・定常直流の近似では電位は一価な関数であり経路に依存しない。\par

\noindent\textbf{2：適切} 真性半導体では温度上昇で電気抵抗が減る。。熱励起により担い手が増え、真性領域では抵抗が減る。\par

\noindent\textbf{3：適切} 純金属に少量の不純物を混ぜると抵抗は増える。。不純物による電子散乱が増え、残留抵抗等が大きくなる。\par

\noindent\textbf{4：不適切} 多くの導体の抵抗率は温度の二乗に比例する。。室温近傍の金属では基準温度からの変化を一次近似することが多い。特定低温域のT\textsuperscript{2}項を多くの導体の一般則にできない。\par

\noindent\textbf{5：適切} 消費電力は単位時間あたり消費するエネルギーである。。P=\allowbreak{}dW/dtであり、単位J/s=\allowbreak{}Wである。\par

\paragraph{再利用する考え方} 抵抗の温度依存と材料の区別を、成立条件・単位・符号から再現する。電場と電束密度、磁場と磁束密度、電荷一定と電圧一定を読み分ける。

\paragraph{共通理論} \hyperref[theory:em:resistivity]{第\ref*{theory:em:resistivity}節}, \hyperref[theory:em:potential]{第\ref*{theory:em:potential}節}

\par 参考：\cite{mit_em06,mit_semiconductor}。

\par 難易度 2（校準初稿）、優先度未評価。

\subsection{令和2年度 問4}
\label{q:r02-reg-applied-04}

\noindent\textbf{原典：} sources/R02/17\_applied\_science.pdf、PDF物理2頁。

\textbf{公式正答：} 1（single）。

\paragraph{問題の要旨} 直列RLC回路のインピーダンスの大きさが最小となる交流周波数を求める。

\Needspace{4\baselineskip}\paragraph{確認した条件・記述（独自要約）}\mbox{}\par

\noindent\textbf{抵抗：} 2 kΩ\par

\noindent\textbf{インダクタンス：} 5 mH\par

\noindent\textbf{容量：} 2 μF\par

\noindent\textbf{接続：} 抵抗・コイル・コンデンサーを直列。理想的な周波数一定素子の定常交流。\par

\paragraph{問われる核心} 直列RLCのリアクタンス相殺

\paragraph{必要知識} 複素インピーダンス、角周波数と周波数、m=\allowbreak{}10\textsuperscript{-3}、μ=10\textsuperscript{-6}。

\paragraph{考え方・正答を選ぶ理由（査読待ち初稿）} \(Z=R+j(\omega L-1/(\omega C))\)なので大きさの二乗は\(R^2+(\omega L-1/(\omega C))^2\)となる。非負の第2項が0となる\(f_0=1/(2\pi\sqrt{LC})\)で最小である。mHとμFをHとFへ換算すると約1.5915 kHz、1を選ぶ。Rは最小値を決めるが、この理想モデルの共振周波数は変えない。

\Needspace{4\baselineskip}\paragraph{選択肢の検討（査読待ち初稿）}\mbox{}\par

\noindent\textbf{1：条件に一致} 1.6 kHz。リアクタンスが相殺する約1.5915 kHzに最も近い。\par

\noindent\textbf{2：条件に不一致} 5.1 kHz。SIへ換算した共振周波数は約1.5915 kHzであり、この候補より離れる。\par

\noindent\textbf{3：条件に不一致} 10 kHz。SIへ換算した共振周波数は約1.5915 kHzであり、この候補より離れる。\par

\noindent\textbf{4：条件に不一致} 100 kHz。SIへ換算した共振周波数は約1.5915 kHzであり、この候補より離れる。\par

\noindent\textbf{5：条件に不一致} 630 kHz。SIへ換算した共振周波数は約1.5915 kHzであり、この候補より離れる。\par

\paragraph{再利用する考え方} 直列RLCのリアクタンス相殺を、成立条件・単位・符号から再現する。電場と電束密度、磁場と磁束密度、電荷一定と電圧一定を読み分ける。

\paragraph{共通理論} \hyperref[theory:em:rlc]{第\ref*{theory:em:rlc}節}

\par 参考：\cite{mit_em12}。

\paragraph{関連問題} \hyperref[q:h29-reg-applied-03]{H29-REG-APPLIED-03}, \hyperref[q:r06-reg-applied-03]{R06-REG-APPLIED-03}

\par 難易度 2（校準初稿）、優先度未評価。

\subsection{令和4年度 問4}
\label{q:r04-reg-applied-04}

\noindent\textbf{原典：} sources/R04/17\_applied\_science.pdf、PDF物理2頁。

\textbf{公式正答：} 3（single）。

\paragraph{問題の要旨} 充電したコンデンサーを抵抗0のコイルへ接続し、最大電流時のコイルエネルギーを求める。

\Needspace{4\baselineskip}\paragraph{確認した条件・記述（独自要約）}\mbox{}\par

\noindent\textbf{容量：} 2.0×10\textsuperscript{-10} F\par

\noindent\textbf{初期電圧：} 5.0 V\par

\noindent\textbf{インダクタンス：} 5.0×10\textsuperscript{-3} H\par

\noindent\textbf{抵抗：} 0。理想LC振動。\par

\paragraph{問われる核心} LC振動の電気・磁気エネルギー交換

\paragraph{必要知識} 無散逸振動、コンデンサーとコイルのエネルギー、最大値の時点。

\paragraph{考え方・正答を選ぶ理由（査読待ち初稿）} 抵抗0なので全エネルギー\(Q^2/(2C)+LI^2/2\)は一定である。電流最大時はQ=\allowbreak{}0、初期の\(CV_0^2/2=(2.0\times10^{-10})5.0^2/2=2.5\times10^{-9}\) Jが全てコイルにある。3を選ぶ。Lは最大電流を決めるが、問われる全エネルギーは初期条件から決まる。

\Needspace{4\baselineskip}\paragraph{選択肢の検討（査読待ち初稿）}\mbox{}\par

\noindent\textbf{1：条件に不一致} 0 J。最大電流時はコイルの磁場エネルギーが最大であり0でない。\par

\noindent\textbf{2：条件に不一致} \(1.3\times10^{-9}\) J。初期エネルギーの約半分だが、最大電流では電気エネルギーは0となり全量が磁気側へ移る。\par

\noindent\textbf{3：条件に一致} \(2.5\times10^{-9}\) J。初期のCV\textsuperscript{2}/2=\allowbreak{}2.5×10\textsuperscript{-9} Jに一致する。\par

\noindent\textbf{4：条件に不一致} \(5.0\times10^{-9}\) J。係数1/2を落としたCV\textsuperscript{2}の値である。\par

\noindent\textbf{5：条件に不一致} \(7.5\times10^{-9}\) J。保存される初期エネルギーの3倍である。\par

\paragraph{再利用する考え方} LC振動の電気・磁気エネルギー交換を、成立条件・単位・符号から再現する。電場と電束密度、磁場と磁束密度、電荷一定と電圧一定を読み分ける。

\paragraph{共通理論} \hyperref[theory:em:lc]{第\ref*{theory:em:lc}節}, \hyperref[theory:em:energy]{第\ref*{theory:em:energy}節}

\par 参考：\cite{mit_em11}。

\par 難易度 2（校準初稿）、優先度未評価。

\subsection{令和6年度 問3}
\label{q:r06-reg-applied-03}

\noindent\textbf{原典：} sources/R06/17\_applied\_science.pdf、PDF物理2頁。

\textbf{公式正答：} 1（single）。

\paragraph{問題の要旨} 直列RLC回路のインピーダンスの大きさが最小となる交流周波数を求める。

\Needspace{4\baselineskip}\paragraph{確認した条件・記述（独自要約）}\mbox{}\par

\noindent\textbf{抵抗：} 200 Ω\par

\noindent\textbf{インダクタンス：} 5.0 mH\par

\noindent\textbf{容量：} 0.50 μF\par

\noindent\textbf{接続：} 抵抗・コイル・コンデンサーを直列。理想的な周波数一定素子の定常交流。\par

\paragraph{問われる核心} 直列RLCのリアクタンス相殺

\paragraph{必要知識} 複素インピーダンス、角周波数と周波数、m=\allowbreak{}10\textsuperscript{-3}、μ=10\textsuperscript{-6}。

\paragraph{考え方・正答を選ぶ理由（査読待ち初稿）} \(Z=R+j(\omega L-1/(\omega C))\)なので大きさの二乗は\(R^2+(\omega L-1/(\omega C))^2\)となる。非負の第2項が0となる\(f_0=1/(2\pi\sqrt{LC})\)で最小である。mHとμFをHとFへ換算すると約3.1831 kHz、1を選ぶ。Rは最小値を決めるが、この理想モデルの共振周波数は変えない。

\Needspace{4\baselineskip}\paragraph{選択肢の検討（査読待ち初稿）}\mbox{}\par

\noindent\textbf{1：条件に一致} 3.2 kHz。リアクタンスが相殺する約3.1831 kHzに最も近い。\par

\noindent\textbf{2：条件に不一致} 6.3 kHz。SIへ換算した共振周波数は約3.1831 kHzであり、この候補より離れる。\par

\noindent\textbf{3：条件に不一致} 20 kHz。SIへ換算した共振周波数は約3.1831 kHzであり、この候補より離れる。\par

\noindent\textbf{4：条件に不一致} 64 kHz。SIへ換算した共振周波数は約3.1831 kHzであり、この候補より離れる。\par

\noindent\textbf{5：条件に不一致} 130 kHz。SIへ換算した共振周波数は約3.1831 kHzであり、この候補より離れる。\par

\paragraph{再利用する考え方} 直列RLCのリアクタンス相殺を、成立条件・単位・符号から再現する。電場と電束密度、磁場と磁束密度、電荷一定と電圧一定を読み分ける。

\paragraph{共通理論} \hyperref[theory:em:rlc]{第\ref*{theory:em:rlc}節}

\par 参考：\cite{mit_em12}。

\paragraph{関連問題} \hyperref[q:r02-reg-applied-04]{R02-REG-APPLIED-04}

\par 難易度 2（校準初稿）、優先度未評価。

\section{自発分極・圧電と強誘電相転移}

\subsection{令和7年度 問7}
\label{q:r07-reg-applied-07}

\noindent\textbf{原典：} sources/R07/17\_applied\_science.pdf、PDF物理3頁。

\textbf{公式正答：} 3（single）。

\par\textbf{検討未確定：} 多ドメイン試料の正味分極、有限外場での飽和分極の測定条件を留保。一次転移の理論的秩序変数の不連続と無条件に同一視しない。

\paragraph{問題の要旨} 自発分極・圧電性と強誘電相転移の包含関係を読む。

\Needspace{4\baselineskip}\paragraph{確認した条件・記述（独自要約）}\mbox{}\par

\noindent\textbf{原典確認：} S7\_051の11原典画像で13問の全五肢・指数・単位・否定を照合。\par

\noindent\textbf{対象：} 自発分極・圧電性と強誘電相転移の包含関係を読む。\par

\paragraph{問われる核心} 自発分極・圧電性と強誘電相転移の包含関係を読む。

\paragraph{必要知識} 3は包含関係を逆にした全称命題。圧電性は応力と分極・電場とひずみの結合、強誘電性は外場なしの分極とその電場反転を特徴とする。磁性のCurie温度と、分極のCurie温度も異なる秩序変数である。

\paragraph{考え方・正答を選ぶ理由（査読待ち初稿）} 3は包含関係を逆にした全称命題。圧電性は応力と分極・電場とひずみの結合、強誘電性は外場なしの分極とその電場反転を特徴とする。磁性のCurie温度と、分極のCurie温度も異なる秩序変数である。

\Needspace{4\baselineskip}\paragraph{選択肢の検討（査読待ち初稿）}\mbox{}\par

\noindent\textbf{1：代表的条件で適切} 外場なしで自発的な双極子。強誘電相の各分極ドメインに自発分極がある。多ドメイン試料で正味の分極が相殺されても局所の自発分極は否定されない。\par

\noindent\textbf{2：代表的条件で適切} 一般にCurie温度以上は常誘電相。典型的な強誘電体の相転移の説明。全物質の全温度経路の普遍的な単一転移へは広げない。\par

\noindent\textbf{3：不適切} 圧電結晶はすべて強誘電体。石英のように圧電性を持ちながら電場で反転する自発分極を持たない結晶がある。強誘電体は圧電性を持つという通常の関係の逆を保証しない。\par

\noindent\textbf{4：代表的条件で適切} 一次転移点で飽和分極が不連続。典型的な一次の強誘電相転移では秩序変数が不連続に変わる。有限電場での測定量や履歴の条件と熱力学の相境界を区別する。\par

\noindent\textbf{5：代表的条件で適切} 強誘電転移は構造相転移。通常の強誘電性ではイオン位置・配向秩序等の構造と分極が関係する。変位型と秩序無秩序型等の微視的機構を区別する。\par

\paragraph{再利用する考え方} 場・応答・担い手・材料・温度を区別し、数値の指数と単位、代表モデルの適用条件を確認する。

\paragraph{共通理論} \hyperref[theory:em:dielectric]{第\ref*{theory:em:dielectric}節}, \hyperref[theory:em:ferroelectric]{第\ref*{theory:em:ferroelectric}節}

\par 参考：\cite{mit_ferroelectric,em51_quartz,em51_ferro_transition}。

\par 難易度 3（校準初稿）、優先度未評価。

\section{磁気模型・測定と熱電輸送則}

\subsection{平成23年度 問2}
\label{q:h23-reg-applied-02}

\noindent\textbf{原典：} sources/H23/17\_applied\_science.pdf、PDF物理1頁。

\textbf{公式正答：} 1（single）。

\paragraph{問題の要旨} 磁性の測定・模型と熱電輸送則の役割を区別する。

\Needspace{4\baselineskip}\paragraph{確認した条件・記述（独自要約）}\mbox{}\par

\noindent\textbf{原典確認：} S7\_051の11原典画像で13問の全五肢・指数・単位・否定を照合。\par

\noindent\textbf{対象：} 磁性の測定・模型と熱電輸送則の役割を区別する。\par

\paragraph{問われる核心} 磁性の測定・模型と熱電輸送則の役割を区別する。

\paragraph{必要知識} 輸送係数の関係である1と、磁気モーメントの模型・温度依存・磁気構造の測定を分ける。Wiedemann–Franz則はκ/(σT)=\allowbreak{}Lで、低温の弾性散乱などの条件でLorenz数を比較する。物質の磁性と全く無関係という排他命題へは変えない。

\paragraph{考え方・正答を選ぶ理由（査読待ち初稿）} 輸送係数の関係である1と、磁気モーメントの模型・温度依存・磁気構造の測定を分ける。Wiedemann–Franz則はκ/(σT)=\allowbreak{}Lで、低温の弾性散乱などの条件でLorenz数を比較する。物質の磁性と全く無関係という排他命題へは変えない。

\Needspace{4\baselineskip}\paragraph{選択肢の検討（査読待ち初稿）}\mbox{}\par

\noindent\textbf{1：不適切} Wiedemann–Franz則。金属の熱伝導率と電気伝導率の関係が中心。磁性との相対比較では最も関係が薄い。磁場の影響を受けないという意味ではない。\par

\noindent\textbf{2：代表的条件で適切} Slater–Pauling曲線。合金の電子数・組成と磁気モーメントの関係を表す。\par

\noindent\textbf{3：代表的条件で適切} Curie則。局在磁気モーメントの常磁性磁化率と温度の関係。\par

\noindent\textbf{4：代表的条件で適切} Ising模型。離散的なスピン変数と相互作用で磁気秩序・相転移を扱う模型。\par

\noindent\textbf{5：代表的条件で適切} 中性子回折。中性子の磁気モーメントによる散乱で磁気秩序を調べられる。結晶構造の回折にも使われる。\par

\paragraph{再利用する考え方} 場・応答・担い手・材料・温度を区別し、数値の指数と単位、代表モデルの適用条件を確認する。

\paragraph{共通理論} \hyperref[theory:em:magnetism]{第\ref*{theory:em:magnetism}節}

\par 参考：\cite{em51_mit_magnetic_readings,em51_wf}。

\par 難易度 2（校準初稿）、優先度未評価。

\section{磁気秩序・Curie/Néelと反磁場}

\subsection{平成27年度 問7}
\label{q:h27-reg-applied-07}

\noindent\textbf{原典：} sources/H27/17\_applied\_science.pdf、PDF物理3頁。

\textbf{公式正答：} 5（single）。

\par\textbf{検討未確定：} 原典の常磁性全体への表現は広すぎる。PauliやVan Vleck等を含む全常磁性に単一の温度法則を適用せず、局在モーメントのCurie領域として説明する。

\paragraph{問題の要旨} 磁気秩序・残留磁化とCurie則の温度依存を区別する。

\Needspace{4\baselineskip}\paragraph{確認した条件・記述（独自要約）}\mbox{}\par

\noindent\textbf{原典確認：} S7\_051の11原典画像で13問の全五肢・指数・単位・否定を照合。\par

\noindent\textbf{対象：} 磁気秩序・残留磁化とCurie則の温度依存を区別する。\par

\paragraph{問われる核心} 磁気秩序・残留磁化とCurie則の温度依存を区別する。

\paragraph{必要知識} 局在モーメントの通常のCurie常磁性でχ=C/Tとなるので5。自発磁化の消失、反磁性の符号、Néel温度、残留磁化の定義を分ける。

\paragraph{考え方・正答を選ぶ理由（査読待ち初稿）} 局在モーメントの通常のCurie常磁性でχ=C/Tとなるので5。自発磁化の消失、反磁性の符号、Néel温度、残留磁化の定義を分ける。

\Needspace{4\baselineskip}\paragraph{選択肢の検討（査読待ち初稿）}\mbox{}\par

\noindent\textbf{1：代表的条件で適切} 昇温で自発磁化が消えるCurie温度。通常の強磁性相は温度上昇で秩序が弱まり、Curie温度で自発磁化を失う。\par

\noindent\textbf{2：代表的条件で適切} 外場に反対向きの反磁性。線形応答の磁化率が負で、誘起磁化が外場と反対向きになる。\par

\noindent\textbf{3：代表的条件で適切} 反強磁性から常磁性へのNéel温度。反平行の長距離秩序が失われる温度。強磁性のCurie温度とは区別する。\par

\noindent\textbf{4：代表的条件で適切} 磁場ゼロで残る残留磁化。磁化履歴のある試料に、外場を取り去っても残る磁化。全強磁性試料が必ず大きな値を持つとは限らない。\par

\noindent\textbf{5：不適切} 常磁性磁化率が温度とともに増加。局在モーメントのCurie領域ではχ=C/Tで、絶対温度の上昇に対して減少する。伝導電子のPauli常磁性の弱い温度依存と混同しない。\par

\paragraph{再利用する考え方} 場・応答・担い手・材料・温度を区別し、数値の指数と単位、代表モデルの適用条件を確認する。

\paragraph{共通理論} \hyperref[theory:em:magnetism]{第\ref*{theory:em:magnetism}節}, \hyperref[theory:em:pauli]{第\ref*{theory:em:pauli}節}

\par 参考：\cite{openstax_magnetism,em51_pauli}。

\par 難易度 3（校準初稿）、優先度未評価。

\subsection{平成29年度 問9}
\label{q:h29-reg-applied-09}

\noindent\textbf{原典：} sources/H29/17\_applied\_science.pdf、PDF物理3頁。

\textbf{公式正答：} 2（single）。

\par\textbf{検討未確定：} 元素の分類は通常の室温・安定相を想定する。常磁性の全モデルへの逆温度則の一般化、黒鉛の一様場と勾配場の混同を避ける。

\paragraph{問題の要旨} Néel温度・反磁場の形状依存・元素の磁性を照合する。

\Needspace{4\baselineskip}\paragraph{確認した条件・記述（独自要約）}\mbox{}\par

\noindent\textbf{原典確認：} S7\_051の11原典画像で13問の全五肢・指数・単位・否定を照合。\par

\noindent\textbf{対象：} Néel温度・反磁場の形状依存・元素の磁性を照合する。\par

\paragraph{問われる核心} Néel温度・反磁場の形状依存・元素の磁性を照合する。

\paragraph{必要知識} 適切な肢2は反強磁性秩序を失うNéel温度。ほかはCurie領域の逆温度依存、MnとCoの違い、形状による反磁場、黒鉛に働き得る磁場勾配力を確認する。

\paragraph{考え方・正答を選ぶ理由（査読待ち初稿）} 適切な肢2は反強磁性秩序を失うNéel温度。ほかはCurie領域の逆温度依存、MnとCoの違い、形状による反磁場、黒鉛に働き得る磁場勾配力を確認する。

\Needspace{4\baselineskip}\paragraph{選択肢の検討（査読待ち初稿）}\mbox{}\par

\noindent\textbf{1：不適切} 常磁性磁化率が温度とともに増加。局在モーメントのCurie領域ではχ=C/Tで、絶対温度の上昇に対して減少する。伝導電子のPauli常磁性の弱い温度依存と混同しない。\par

\noindent\textbf{2：代表的条件で適切} 反強磁性から常磁性へのNéel温度。反平行の長距離秩序が失われる温度。強磁性のCurie温度とは区別する。\par

\noindent\textbf{3：不適切} Fe・Mn・Niが単体強磁性。通常の室温の代表的強磁性元素はFe・Co・Ni。原典のMnをCoと読み替えない。Mnの相・温度依存は別に扱う。\par

\noindent\textbf{4：不適切} 反磁場が形状に依存しない。磁化で生じる反磁場は試料形状と磁化分布に依存する。均一磁化した楕円体でHd=−NMと表せる。\par

\noindent\textbf{5：不適切} 黒鉛は磁場中で全く力を受けない。黒鉛は反磁性を持ち、不均一な磁場で力を受け得る。一様な理想場では正味の並進力がゼロでも、全条件で力がないという記述は誤り。\par

\paragraph{再利用する考え方} 場・応答・担い手・材料・温度を区別し、数値の指数と単位、代表モデルの適用条件を確認する。

\paragraph{共通理論} \hyperref[theory:em:magnetism]{第\ref*{theory:em:magnetism}節}, \hyperref[theory:em:demag]{第\ref*{theory:em:demag}節}

\par 参考：\cite{openstax_magnetism,em51_pauli,em51_demag}。

\paragraph{関連問題} \hyperref[q:r06-reg-applied-07]{R06-REG-APPLIED-07}

\par 難易度 3（校準初稿）、優先度未評価。

\subsection{令和6年度 問7}
\label{q:r06-reg-applied-07}

\noindent\textbf{原典：} sources/R06/17\_applied\_science.pdf、PDF物理3頁。

\textbf{公式正答：} 3（single）。

\par\textbf{検討未確定：} 元素の分類は通常の室温・安定相を想定する。常磁性の全モデルへの逆温度則の一般化、黒鉛の一様場と勾配場の混同を避ける。

\paragraph{問題の要旨} Néel温度・反磁場の形状依存・元素の磁性を照合する。

\Needspace{4\baselineskip}\paragraph{確認した条件・記述（独自要約）}\mbox{}\par

\noindent\textbf{原典確認：} S7\_051の11原典画像で13問の全五肢・指数・単位・否定を照合。\par

\noindent\textbf{対象：} Néel温度・反磁場の形状依存・元素の磁性を照合する。\par

\paragraph{問われる核心} Néel温度・反磁場の形状依存・元素の磁性を照合する。

\paragraph{必要知識} 適切な肢3は反強磁性秩序を失うNéel温度。ほかはCurie領域の逆温度依存、MnとCoの違い、形状による反磁場、黒鉛に働き得る磁場勾配力を確認する。

\paragraph{考え方・正答を選ぶ理由（査読待ち初稿）} 適切な肢3は反強磁性秩序を失うNéel温度。ほかはCurie領域の逆温度依存、MnとCoの違い、形状による反磁場、黒鉛に働き得る磁場勾配力を確認する。

\Needspace{4\baselineskip}\paragraph{選択肢の検討（査読待ち初稿）}\mbox{}\par

\noindent\textbf{1：不適切} 常磁性磁化率が温度とともに増加。局在モーメントのCurie領域ではχ=C/Tで、絶対温度の上昇に対して減少する。伝導電子のPauli常磁性の弱い温度依存と混同しない。\par

\noindent\textbf{2：不適切} 黒鉛は磁場中で全く力を受けない。黒鉛は反磁性を持ち、不均一な磁場で力を受け得る。一様な理想場では正味の並進力がゼロでも、全条件で力がないという記述は誤り。\par

\noindent\textbf{3：代表的条件で適切} 反強磁性から常磁性へのNéel温度。反平行の長距離秩序が失われる温度。強磁性のCurie温度とは区別する。\par

\noindent\textbf{4：不適切} Fe・Mn・Niが単体強磁性。通常の室温の代表的強磁性元素はFe・Co・Ni。原典のMnをCoと読み替えない。Mnの相・温度依存は別に扱う。\par

\noindent\textbf{5：不適切} 反磁場が形状に依存しない。磁化で生じる反磁場は試料形状と磁化分布に依存する。均一磁化した楕円体でHd=−NMと表せる。\par

\paragraph{再利用する考え方} 場・応答・担い手・材料・温度を区別し、数値の指数と単位、代表モデルの適用条件を確認する。

\paragraph{共通理論} \hyperref[theory:em:magnetism]{第\ref*{theory:em:magnetism}節}, \hyperref[theory:em:demag]{第\ref*{theory:em:demag}節}

\par 参考：\cite{openstax_magnetism,em51_pauli,em51_demag}。

\paragraph{関連問題} \hyperref[q:h29-reg-applied-09]{H29-REG-APPLIED-09}

\par 難易度 3（校準初稿）、優先度未評価。

\section{局在スピン・Fermi統計とPauli応答}

\subsection{令和元年度 問8}
\label{q:r01-reg-applied-08}

\noindent\textbf{原典：} sources/R01/17\_applied\_science.pdf、PDF物理3頁。

\textbf{公式正答：} 5（single）。

\par\textbf{検討未確定：} 肢2の完全反磁性はMeissner状態の条件付き。II型超伝導の混合状態を含めて全状態へ断定しない。

\paragraph{問題の要旨} Pauliのスピン応答と軌道反磁性を区別する。

\Needspace{4\baselineskip}\paragraph{確認した条件・記述（独自要約）}\mbox{}\par

\noindent\textbf{原典確認：} S7\_051の11原典画像で13問の全五肢・指数・単位・否定を照合。\par

\noindent\textbf{対象：} Pauliのスピン応答と軌道反磁性を区別する。\par

\paragraph{問われる核心} Pauliのスピン応答と軌道反磁性を区別する。

\paragraph{必要知識} 誤りは5の起源の取り違え。局在スピンのCurie則と、Fermi準位近傍の伝導電子によるPauli常磁性を分ける。Meissner効果も超伝導の全磁場領域での完全排除へは広げない。

\paragraph{考え方・正答を選ぶ理由（査読待ち初稿）} 誤りは5の起源の取り違え。局在スピンのCurie則と、Fermi準位近傍の伝導電子によるPauli常磁性を分ける。Meissner効果も超伝導の全磁場領域での完全排除へは広げない。

\Needspace{4\baselineskip}\paragraph{選択肢の検討（査読待ち初稿）}\mbox{}\par

\noindent\textbf{1：代表的条件で適切} 室温の鉄は強磁性。通常圧・通常の純Feは室温で強磁性。結晶構造転移と磁気秩序の転移を区別する。\par

\noindent\textbf{2：条件・語義を留保} Meissner効果による完全反磁性。Meissner状態では磁束を排除する。臨界磁場・温度以下の条件が必要で、II型の混合状態では量子化磁束が侵入する。\par

\noindent\textbf{3：代表的条件で適切} 反強磁性から常磁性へのNéel温度。反平行の長距離秩序が失われる温度。強磁性のCurie温度とは区別する。\par

\noindent\textbf{4：代表的条件で適切} Curie点以上で常磁性。通常の強磁性物質ではCurie温度を超えると自発磁化秩序を失う。\par

\noindent\textbf{5：不適切} Pauli常磁性が原子核まわりの軌道運動。伝導電子のスピンに外場が作用し、Fermi統計を反映した応答となる。自由電子の軌道応答はLandau反磁性として区別する。\par

\paragraph{再利用する考え方} 場・応答・担い手・材料・温度を区別し、数値の指数と単位、代表モデルの適用条件を確認する。

\paragraph{共通理論} \hyperref[theory:em:pauli]{第\ref*{theory:em:pauli}節}

\par 参考：\cite{em51_pauli,openstax_magnetism}。

\par 難易度 3（校準初稿）、優先度未評価。

\section{金属抵抗率・温度と単位桁}

\subsection{平成26年度 問5}
\label{q:h26-reg-applied-05}

\noindent\textbf{原典：} sources/H26/17\_applied\_science.pdf、PDF物理2頁。

\textbf{公式正答：} 5（single）。

\paragraph{問題の要旨} 金属抵抗の形状・温度依存と抵抗率の桁を照合する。

\Needspace{4\baselineskip}\paragraph{確認した条件・記述（独自要約）}\mbox{}\par

\noindent\textbf{原典確認：} S7\_051の11原典画像で13問の全五肢・指数・単位・否定を照合。\par

\noindent\textbf{対象：} 金属抵抗の形状・温度依存と抵抗率の桁を照合する。\par

\paragraph{問われる核心} 金属抵抗の形状・温度依存と抵抗率の桁を照合する。

\paragraph{必要知識} 各回の誤りは肢5の抵抗率の桁である。R=ρL/AとI=\allowbreak{}V/Rを、材料・断面・温度一定という条件から使う。残留抵抗と温度依存を分け、H26のΩcm、R01とR05のΩmを独立に校正した。通常金属の温度傾向を特殊な転移や全合金へ広げない。

\paragraph{考え方・正答を選ぶ理由（査読待ち初稿）} 各回の誤りは肢5の抵抗率の桁である。R=ρL/AとI=\allowbreak{}V/Rを、材料・断面・温度一定という条件から使う。残留抵抗と温度依存を分け、H26のΩcm、R01とR05のΩmを独立に校正した。通常金属の温度傾向を特殊な転移や全合金へ広げない。

\Needspace{4\baselineskip}\paragraph{選択肢の検討（査読待ち初稿）}\mbox{}\par

\noindent\textbf{1：代表的条件で適切} 300Kと100Kの抵抗。通常の金属では格子振動による散乱が高温ほど増え、室温側の抵抗が大きい。相転移や異常な材料を含む全温度の保証ではない。\par

\noindent\textbf{2：代表的条件で適切} 不純物・格子欠陥と残留抵抗。低温で格子振動の寄与が減っても不純物・欠陥による散乱は残る。組成・構造による寄与を考える。\par

\noindent\textbf{3：代表的条件で適切} 一定断面でRが長さに比例。均一な導線はR=ρL/A。同じ材料・温度・断面積を保つ条件が必要。\par

\noindent\textbf{4：代表的条件で適切} 電流が電位差に比例。オーム性・温度一定ならI=\allowbreak{}V/R。自己加熱でRが変わる高電圧域まで無条件には広げない。\par

\noindent\textbf{5：不適切} 室温の金属が10\textsuperscript{2}〜10\textsuperscript{4}Ωcm。原典は正の指数でΩcm。Cuの代表値1.68×10\textsuperscript{-8}Ωmは1.68×10\textsuperscript{-6}Ωcm。元の範囲は1〜100Ωmに相当し、金属の典型抵抗率から大きく外れる。\par

\paragraph{再利用する考え方} 場・応答・担い手・材料・温度を区別し、数値の指数と単位、代表モデルの適用条件を確認する。

\paragraph{共通理論} \hyperref[theory:em:resistivity]{第\ref*{theory:em:resistivity}節}

\par 参考：\cite{openstax_resistivity,em51_resistance}。

\paragraph{関連問題} \hyperref[q:r01-reg-applied-07]{R01-REG-APPLIED-07}

\par 難易度 2（校準初稿）、優先度未評価。

\subsection{令和元年度 問7}
\label{q:r01-reg-applied-07}

\noindent\textbf{原典：} sources/R01/17\_applied\_science.pdf、PDF物理3頁。

\textbf{公式正答：} 3（single）。

\paragraph{問題の要旨} 金属抵抗の形状・温度依存と抵抗率の桁を照合する。

\Needspace{4\baselineskip}\paragraph{確認した条件・記述（独自要約）}\mbox{}\par

\noindent\textbf{原典確認：} S7\_051の11原典画像で13問の全五肢・指数・単位・否定を照合。\par

\noindent\textbf{対象：} 金属抵抗の形状・温度依存と抵抗率の桁を照合する。\par

\paragraph{問われる核心} 金属抵抗の形状・温度依存と抵抗率の桁を照合する。

\paragraph{必要知識} 各回の誤りは肢3の抵抗率の桁である。R=ρL/AとI=\allowbreak{}V/Rを、材料・断面・温度一定という条件から使う。残留抵抗と温度依存を分け、H26のΩcm、R01とR05のΩmを独立に校正した。通常金属の温度傾向を特殊な転移や全合金へ広げない。

\paragraph{考え方・正答を選ぶ理由（査読待ち初稿）} 各回の誤りは肢3の抵抗率の桁である。R=ρL/AとI=\allowbreak{}V/Rを、材料・断面・温度一定という条件から使う。残留抵抗と温度依存を分け、H26のΩcm、R01とR05のΩmを独立に校正した。通常金属の温度傾向を特殊な転移や全合金へ広げない。

\Needspace{4\baselineskip}\paragraph{選択肢の検討（査読待ち初稿）}\mbox{}\par

\noindent\textbf{1：代表的条件で適切} 一定断面でRが長さに比例。均一な導線はR=ρL/A。同じ材料・温度・断面積を保つ条件が必要。\par

\noindent\textbf{2：代表的条件で適切} 電流が電位差に比例。オーム性・温度一定ならI=\allowbreak{}V/R。自己加熱でRが変わる高電圧域まで無条件には広げない。\par

\noindent\textbf{3：不適切} 室温の金属が10\textsuperscript{2}〜10\textsuperscript{4}Ωm。原典はΩmで、H26のΩcmと同じ数値扱いにはしない。一般的金属の代表値10\textsuperscript{-8}〜10\textsuperscript{-6}Ωm程度と桁が違う。\par

\noindent\textbf{4：代表的条件で適切} 300Kと100Kの抵抗。通常の金属では格子振動による散乱が高温ほど増え、室温側の抵抗が大きい。相転移や異常な材料を含む全温度の保証ではない。\par

\noindent\textbf{5：代表的条件で適切} 不純物・格子欠陥と残留抵抗。低温で格子振動の寄与が減っても不純物・欠陥による散乱は残る。組成・構造による寄与を考える。\par

\paragraph{再利用する考え方} 場・応答・担い手・材料・温度を区別し、数値の指数と単位、代表モデルの適用条件を確認する。

\paragraph{共通理論} \hyperref[theory:em:resistivity]{第\ref*{theory:em:resistivity}節}

\par 参考：\cite{openstax_resistivity,em51_resistance}。

\paragraph{関連問題} \hyperref[q:h26-reg-applied-05]{H26-REG-APPLIED-05}, \hyperref[q:r05-reg-applied-08]{R05-REG-APPLIED-08}

\par 難易度 2（校準初稿）、優先度未評価。

\subsection{令和5年度 問8}
\label{q:r05-reg-applied-08}

\noindent\textbf{原典：} sources/R05/17\_applied\_science.pdf、PDF物理3頁。

\textbf{公式正答：} 4（single）。

\paragraph{問題の要旨} 金属抵抗の形状・温度依存と抵抗率の桁を照合する。

\Needspace{4\baselineskip}\paragraph{確認した条件・記述（独自要約）}\mbox{}\par

\noindent\textbf{原典確認：} S7\_051の11原典画像で13問の全五肢・指数・単位・否定を照合。\par

\noindent\textbf{対象：} 金属抵抗の形状・温度依存と抵抗率の桁を照合する。\par

\paragraph{問われる核心} 金属抵抗の形状・温度依存と抵抗率の桁を照合する。

\paragraph{必要知識} 各回の誤りは肢4の抵抗率の桁である。R=ρL/AとI=\allowbreak{}V/Rを、材料・断面・温度一定という条件から使う。残留抵抗と温度依存を分け、H26のΩcm、R01とR05のΩmを独立に校正した。通常金属の温度傾向を特殊な転移や全合金へ広げない。

\paragraph{考え方・正答を選ぶ理由（査読待ち初稿）} 各回の誤りは肢4の抵抗率の桁である。R=ρL/AとI=\allowbreak{}V/Rを、材料・断面・温度一定という条件から使う。残留抵抗と温度依存を分け、H26のΩcm、R01とR05のΩmを独立に校正した。通常金属の温度傾向を特殊な転移や全合金へ広げない。

\Needspace{4\baselineskip}\paragraph{選択肢の検討（査読待ち初稿）}\mbox{}\par

\noindent\textbf{1：代表的条件で適切} 不純物・格子欠陥と残留抵抗。低温で格子振動の寄与が減っても不純物・欠陥による散乱は残る。組成・構造による寄与を考える。\par

\noindent\textbf{2：代表的条件で適切} 一定断面でRが長さに比例。均一な導線はR=ρL/A。同じ材料・温度・断面積を保つ条件が必要。\par

\noindent\textbf{3：代表的条件で適切} −195℃と0℃の抵抗。それぞれ約78Kと273K。通常の金属の比較では0℃側が大きい。摂氏温度を比で使わない。\par

\noindent\textbf{4：不適切} 室温の金属が1〜10\textsuperscript{2}Ωm。原典は1〜100Ωm。一般的な金属の抵抗率の典型値より何桁も大きい。特定合金の値を全金属へ固定しない。\par

\noindent\textbf{5：代表的条件で適切} 電流が電位差に比例。オーム性・温度一定ならI=\allowbreak{}V/R。自己加熱でRが変わる高電圧域まで無条件には広げない。\par

\paragraph{再利用する考え方} 場・応答・担い手・材料・温度を区別し、数値の指数と単位、代表モデルの適用条件を確認する。

\paragraph{共通理論} \hyperref[theory:em:resistivity]{第\ref*{theory:em:resistivity}節}

\par 参考：\cite{openstax_resistivity,em51_resistance}。

\paragraph{関連問題} \hyperref[q:r01-reg-applied-07]{R01-REG-APPLIED-07}

\par 難易度 2（校準初稿）、優先度未評価。

\section{誘電損失・渦電流と表皮深さ}

\subsection{令和3年度 問3}
\label{q:r03-reg-applied-03}

\noindent\textbf{原典：} sources/R03/17\_applied\_science.pdf、PDF物理1頁。

\textbf{公式正答：} 1（single）。

\par\textbf{検討未確定：} 公式1を保持するが肢3の全導電体への履歴損、肢4の厳密な侵入ゼロには疑義。無電極トーチの原回路も未照合でAstraへ残す。

\paragraph{問題の要旨} 誘電損失・渦電流と表皮深さを分けて高周波加熱を読む。

\Needspace{4\baselineskip}\paragraph{確認した条件・記述（独自要約）}\mbox{}\par

\noindent\textbf{原典確認：} S7\_051の11原典画像で13問の全五肢・指数・単位・否定を照合。\par

\noindent\textbf{対象：} 誘電損失・渦電流と表皮深さを分けて高周波加熱を読む。\par

\paragraph{問われる核心} 誘電損失・渦電流と表皮深さを分けて高周波加熱を読む。

\paragraph{必要知識} 公式1は誘電体と導電体の入れ替えに対応する。誘電損失、導電電流のJoule損、磁化履歴損を別々に考える。表皮深さδが有限であるため、肢4の内部へ全く侵入しないという字義を正確な物理定義として教えない。

\paragraph{考え方・正答を選ぶ理由（査読待ち初稿）} 公式1は誘電体と導電体の入れ替えに対応する。誘電損失、導電電流のJoule損、磁化履歴損を別々に考える。表皮深さδが有限であるため、肢4の内部へ全く侵入しないという字義を正確な物理定義として教えない。

\Needspace{4\baselineskip}\paragraph{選択肢の検討（査読待ち初稿）}\mbox{}\par

\noindent\textbf{1：不適切} 誘電加熱の対象を導電体とする。代表的な誘電加熱は誘電体の交流電場に対する損失を利用する。導電体を主対象とする誘導加熱と入れ替えている。有限導電率の試料で両方の損失が共存することはある。\par

\noindent\textbf{2：代表的条件で適切} GHz帯の誘電加熱をマイクロ波加熱。マイクロ波電場で誘電損失を生じさせる方式。高周波誘電加熱の全周波数がGHzという意味ではない。\par

\noindent\textbf{3：条件・語義を留保} 誘導加熱で渦電流損・履歴損。導電体に誘起された渦電流のJoule損が基本。磁化履歴による損失は強磁性体等で寄与し、すべての導電体が履歴損を持つわけではない。\par

\noindent\textbf{4：条件・語義を留保} 内部に侵入しない表皮効果。電流・場は表面近傍へ集中するが、良導体半空間モデルでも深さに対し指数的に減衰し、数学的にゼロの侵入ではない。原典の近似表現を留保する。\par

\noindent\textbf{5：代表的条件で適切} 無極放電・融合プラズマの高周波加熱。プラズマにも高周波電磁場からエネルギーを供給できる。QST資料で融合プラズマのRF加熱を確認。トーチの具体的な無電極回路は別途照合する。\par

\paragraph{再利用する考え方} 場・応答・担い手・材料・温度を区別し、数値の指数と単位、代表モデルの適用条件を確認する。

\paragraph{共通理論} \hyperref[theory:em:heating]{第\ref*{theory:em:heating}節}

\par 参考：\cite{em51_heating,em51_skin,em51_plasma}。

\par 難易度 3（校準初稿）、優先度未評価。

\section{理想変圧器の電圧巻数比}

\subsection{令和5年度 問4}
\label{q:r05-reg-applied-04}

\noindent\textbf{原典：} sources/R05/17\_applied\_science.pdf、PDF物理2頁。

\textbf{公式正答：} 3（single）。

\paragraph{問題の要旨} RLC共振・ソレノイド・理想変圧器の巻数比を照合する。

\Needspace{4\baselineskip}\paragraph{確認した条件・記述（独自要約）}\mbox{}\par

\noindent\textbf{原典確認：} S7\_051の11原典画像で13問の全五肢・指数・単位・否定を照合。\par

\noindent\textbf{対象：} RLC共振・ソレノイド・理想変圧器の巻数比を照合する。\par

\paragraph{問われる核心} RLC共振・ソレノイド・理想変圧器の巻数比を照合する。

\paragraph{必要知識} 誤りは3の逆比。各巻の起電力が共通の磁束変化に比例するので、二次の巻数が多ければ二次電圧も大きい。共振の駆動条件と長いソレノイドの近似条件も別々に確認する。

\paragraph{考え方・正答を選ぶ理由（査読待ち初稿）} 誤りは3の逆比。各巻の起電力が共通の磁束変化に比例するので、二次の巻数が多ければ二次電圧も大きい。共振の駆動条件と長いソレノイドの近似条件も別々に確認する。

\Needspace{4\baselineskip}\paragraph{選択肢の検討（査読待ち初稿）}\mbox{}\par

\noindent\textbf{1：代表的条件で適切} 直列RLCのf=\allowbreak{}1/(2π√LC)。交流電源で駆動する直列回路の電流共振ではωL=\allowbreak{}1/(ωC)。抵抗で減衰する自由振動周波数とは別である。\par

\noindent\textbf{2：代表的条件で適切} ソレノイドBが巻数密度nに比例。長く端効果が小さい、電流・透磁率一定のモデルでB=μnI。飽和域でμを一定としない。\par

\noindent\textbf{3：不適切} E2=\allowbreak{}(N1/N2)E1。理想変圧器は同じ時間変化磁束を共有し、E2/E1=\allowbreak{}N2/N1。原典は分子分母が逆。電圧の最大値同士でも同じ巻数比になる。\par

\noindent\textbf{4：代表的条件で適切} 弱い磁化が外場と同じ向きなら常磁性。通常の線形応答で小さな正の磁化率を持つ。外場なしの自発磁化を持つ強磁性と分ける。\par

\noindent\textbf{5：代表的条件で適切} Fe・Ni・Co等の強磁性体をコアに利用。代表的な強磁性元素で磁束を集中できる。実際の交流コアの選定では透磁率だけでなく保磁力・渦電流損・周波数を考える。\par

\paragraph{再利用する考え方} 場・応答・担い手・材料・温度を区別し、数値の指数と単位、代表モデルの適用条件を確認する。

\paragraph{共通理論} \hyperref[theory:em:rlc]{第\ref*{theory:em:rlc}節}, \hyperref[theory:em:solenoid]{第\ref*{theory:em:solenoid}節}, \hyperref[theory:em:transformer]{第\ref*{theory:em:transformer}節}

\par 参考：\cite{mit_em12,mit_em09,em51_transformer_ref,openstax_magnetism}。

\par 難易度 2（校準初稿）、優先度未評価。

\section{熱電効果・可逆熱と散逸}

\subsection{平成29年度 問7}
\label{q:h29-reg-applied-07}

\noindent\textbf{原典：} sources/H29/17\_applied\_science.pdf、PDF物理3頁。

\textbf{公式正答：} 5（single）。

\par\textbf{検討未確定：} 候補4は材料・端点温度・開回路を省略している。起電力と負荷時の端子電圧の区別を付し、一般に形状が無関係とはしない。

\paragraph{問題の要旨} ゼーベック・ペルティエ・トムソン効果を、温度依存と電流依存で区別する。

\Needspace{4\baselineskip}\paragraph{確認した条件・記述（独自要約）}\mbox{}\par

\noindent\textbf{確認範囲：} 問題幹、数値、式、否定、全5候補を原典画像で照合した。\par

\noindent\textbf{成立条件：} 熱起電力は開回路で同じ材料・同じ両端温度を比較する。端点温度と材料を固定した熱起電力と、負荷を接続した端子電圧を区別する。\par

\paragraph{問われる核心} ゼーベック・ペルティエ・トムソン効果を、温度依存と電流依存で区別する。

\paragraph{必要知識} 熱起電力は開回路で同じ材料・同じ両端温度を比較する。端点温度と材料を固定した熱起電力と、負荷を接続した端子電圧を区別する。

\paragraph{考え方・正答を選ぶ理由（査読待ち初稿）} 温度勾配を持つ導体のトムソン熱は電流に一次であり、電流反転で吸熱と発熱が入れ替わる。\(I^2R\) に従い電流反転で符号が変わらないのはジュール熱である。したがって5が不適切である。開回路の熱起電力は材料と端点温度が同じなら熱電能の温度積分で決まり、長さ・断面積を含まない。

\Needspace{4\baselineskip}\paragraph{選択肢の検討（査読待ち初稿）}\mbox{}\par

\noindent\textbf{1：適切} ゼーベック効果は温度差測定に利用。異種材料の接点温度差から生じる起電力を熱電対の測温に用いる。\par

\noindent\textbf{2：適切} ペルティエ効果は熱電冷却に応用。異種材料の接合部で電流方向に応じて吸熱又は発熱が起こる。\par

\noindent\textbf{3：適切} 熱電能は温度依存。ゼーベック係数S(T)は一般に材料と温度に依存し、一定値とするのは限定した温度範囲の近似である。\par

\noindent\textbf{4：条件付きで適切} 熱起電力は長さや太さに無関係。材料と端点温度を固定した開回路の起電力には幾何形状が直接入らない。電気抵抗や負荷時の端子電圧、温度分布まで形状に依存しないという意味ではない。\par

\noindent\textbf{5：不適切} トムソン熱は電流の2乗に比例。トムソン熱は電流と温度勾配に一次で可逆である。二乗依存はジュール熱の特徴である。\par

\paragraph{再利用する考え方} 接合部のペルティエ熱、単一導体内のトムソン熱、散逸によるジュール熱の場所と電流反転時の符号を比べる。

\paragraph{共通理論} \hyperref[theory:electrochem:seebeck]{第\ref*{theory:electrochem:seebeck}節}, \hyperref[theory:electrochem:peltier]{第\ref*{theory:electrochem:peltier}節}, \hyperref[theory:electrochem:thomson]{第\ref*{theory:electrochem:thomson}節}, \hyperref[theory:electrochem:joule]{第\ref*{theory:electrochem:joule}節}

\par 参考：\cite{mit_thermoelectric2024}。

\paragraph{関連問題} \hyperref[q:r03-reg-applied-06]{R03-REG-APPLIED-06}

\par 難易度 2（校準初稿）、優先度未評価。

\subsection{令和3年度 問6}
\label{q:r03-reg-applied-06}

\noindent\textbf{原典：} sources/R03/17\_applied\_science.pdf、PDF物理3頁。

\textbf{公式正答：} 2（single）。

\paragraph{問題の要旨} 熱電現象の名称から光学現象を区別する。

\Needspace{4\baselineskip}\paragraph{確認した条件・記述（独自要約）}\mbox{}\par

\noindent\textbf{確認範囲：} 問題幹、数値、式、否定、全5候補を原典画像で照合した。\par

\noindent\textbf{成立条件：} 最も不適切な用語を選ぶ。カー効果には電気光学と磁気光学の用法があるが、いずれもこの熱電効果の名称ではない。\par

\paragraph{問われる核心} 熱電現象の名称から光学現象を区別する。

\paragraph{必要知識} 最も不適切な用語を選ぶ。カー効果には電気光学と磁気光学の用法があるが、いずれもこの熱電効果の名称ではない。

\paragraph{考え方・正答を選ぶ理由（査読待ち初稿）} ゼーベック、ペルティエ、トムソンと熱起電力はいずれも熱と電気の結合に関わる。カー効果は電場による複屈折又は磁気光学的な偏光変化などの光学現象の名称であり、列挙された熱電現象と区別して2を選ぶ。

\Needspace{4\baselineskip}\paragraph{選択肢の検討（査読待ち初稿）}\mbox{}\par

\noindent\textbf{1：適切} ゼーベック効果。温度差により起電力が生じる熱電現象である。\par

\noindent\textbf{2：不適切} カー効果。光の屈折・偏光に関わる名称であり、温度差と電流の相互作用を表す熱電効果ではない。\par

\noindent\textbf{3：適切} ペルティエ効果。異種材料接合部を流れる電流に伴う吸放熱である。\par

\noindent\textbf{4：適切} トムソン効果。温度勾配を持つ材料内の電流に伴う可逆的な吸放熱である。\par

\noindent\textbf{5：適切} 熱起電力。温度差による起電力そのものを表す用語である。\par

\paragraph{再利用する考え方} 似た人名の効果を、外部の原因と観測量で分類する。熱電は温度・電流・熱流、カー効果は電磁場と光学応答である。

\paragraph{共通理論} \hyperref[theory:electrochem:seebeck]{第\ref*{theory:electrochem:seebeck}節}, \hyperref[theory:electrochem:peltier]{第\ref*{theory:electrochem:peltier}節}, \hyperref[theory:electrochem:thomson]{第\ref*{theory:electrochem:thomson}節}, \hyperref[theory:electrochem:kerr]{第\ref*{theory:electrochem:kerr}節}

\par 参考：\cite{mit_thermoelectric2024,mit_kerr_calendar2005}。

\paragraph{関連問題} \hyperref[q:h29-reg-applied-07]{H29-REG-APPLIED-07}

\par 難易度 2（校準初稿）、優先度未評価。

